Motivation : modéliser des phénomènes aléatoires a l'aide de 3 outils : 

\uindent{Univers $\Omega$ }

Ensemble de tous les résultats possibles d'une expérience aléatoire.

\uindent{Tribu $\mathcal A \subset \mathscr P(\Omega)$ }

Ensemble de toutes les propriétés "raisonnables" que peut posséder 
$\omega \in \Omega$. 

Un élément de $\mathcal A$ est un évènement. 

$A \in \mathcal A$ est un évènement tel que $A = \{ \omega \in \Omega, \omega \text{ satisfait à l'évènement } \}$


\uindent{Probabilité $P$ }

$\fonction P {\mathcal A}{\ninter 0 1} A {P(A)}$

$P(A)$ évalue la vraisemblance de l'évènement $A$. 

\section{Espaces probabilisé}

Soit $\Omega$ un ensemble non vide appelé univers. 

\subsection{Tribu }

\begin{definition}{Tribu}{}
    Soit $\mathcal A \subset \mathscr P(\Omega)$. On dit que $\mathcal A$ est une tribu (ou une $\sigma$-algèbre)
    ssi 
    \begin{enumerate}
        \item $\emptyset \in \mathcal A$
        \item $\forall A \in \mathcal A$, $\Omega \setminus A = \bar A \in \mathcal A$ et $\mathcal A$ est stable 
        par complémentaire ($\bar A = \{ \omega \in \Omega \mid \omega \notin A \}$).
        \item $\mathcal A$ est stable par union dénombrable, c'est-à-dire si $(A_i)_{i \in \N}$ est une suite 
        de $\mathcal A$, alors $\displaystyle \bigcup_{i \in \N} A_i = \{\omega \in \Omega \mid \exists i \in \N, \, \omega \in A_i\} \in \mathcal A$ 
    \end{enumerate}
    On a alors : 
    \begin{enumerate}
        \item $\Omega \in \mathcal A$
        \item Si $(A_i)_{i \in \N} $ suite de $\mathcal A$ alors $\bigcap_{i \in \N} A_i = \{\omega \in \Omega \mid \forall i \in \N, \, \omega \in A_i\} \in \mathcal A$
        \item $\mathcal A$ est stable par union et intersection finie. 
    \end{enumerate}
\end{definition}

\begin{proof}
    Par (1), $\emptyset \in \mathcal A$ donc par (2) $\bar \emptyset = \Omega \in A$

    \avspace Soit $(A_i)_{i \in \N}$ dans $\displaystyle \mathcal A$, on a $\bigcap_{i \in \N} A_i = \overline{\bigcup_{i \in \N} \bar A_i}$

    Or par (2), $\forall i, \, \bar A_i \in \mathcal A$, par (3) $\displaystyle \bigcup_{i \in \N} \bar A_i \in \mathcal A$
    et par (2), $\displaystyle \bigcap_{i \in \N} A_i \in \mathcal A$

    \avspace Soit $\displaystyle (A_i)_{1 \leqslant i \leqslant N}$ famille finie de $\mathcal A$ en posant 
    $\forall i \geqslant N+ 1 $, $A_i = \emptyset \in \mathcal A$ et donc $\displaystyle \bigcup_{i \in \N} A_i \in \mathcal A$
    donc $\displaystyle \bigcup_{i = 1}^N A_i \in \mathcal A$

    \avspace Idem avec l'intersection en posant $A_i = \Omega$
\end{proof}


\begin{exemple}
    $\mathcal A = \{\emptyset, \Omega\}$ est une tribu. 

    \avspace Si $A_0 \in \mathscr P(\Omega)$ telle que $A_0 \neq \emptyset$ et $A_0 \neq \Omega$

    \avspace $\mathcal A = \{\emptyset, A_0, \bar A_0, \Omega\}$ est une tribu, on l'appelle tribu 
    engendrée par $A_0$. 

    \avspace $\mathcal A = \mathscr P(\Omega)$ est une tribu. 

    Si $\Omega$ est fini ou dénombrable c'est celle que l'on choisira. 

    Si $\Omega$ est infini non dénombrable, $\mathscr P(\Omega)$ n'est pas une tribu intéressante. 

    Si $\Omega = \R$ on utilise en général la tribu des boréliens qui est engendrée par les intervalles ouverts. 
\end{exemple}

\begin{definition}{}{}
    Soit $\mathcal A$ une tribu de $\Omega$. 
    \begin{itemize}
        \item $(\Omega, \mathcal A)$ est appelé espace probabilisable. 
        \item Un élément $A$ de $\mathcal A$ est appelé un évènement. 
        \begin{itemize}
            \item $A \subset \Omega$
            \item $A = \{ \omega \in \Omega, \, \omega \in A\}$
        \end{itemize}
        \item $\bar A$ est appelé évènement contraire de $\mathcal A$
        \item $\Omega$ est appelé évènement certain
        \item $\emptyset$ est appelé évènement impossible
    \end{itemize}
\end{definition}

Attention, il y a un problème dans cette définition, $A = \{ \omega \in \Omega, \, \omega \in A\}$ 
n'a aucun sens (à montrer à Quibel un jour).

\subsection{Probabilité}

\begin{definition}{Espace probabilisé}{}
    Soit $(\Omega, \mathcal A)$ un espace probabilisable. On appelle probabilité 
    sur cet espace toute application $\fonction{P}{\mathcal A}{[0, 1]}{A}{P(A)}$
    vérifiant :
    \begin{itemize}
        \item $P(\Omega) = 1$
        \item Pour toute suite $(A_n)$ d'évènements 2 à 2 disjoints, 
        $\displaystyle P\left( \bigcup_{n \in \N} A_n \right) = \sum_{n = 0 }^{\infty    } P(A_n ) = \sum_{n \in \N    } P(A_n)$
        appelé $\sigma$-additivité.
    \end{itemize}

    On dit alors que $(\Omega, \mathcal A, P)$ est un espace probabilisé. 
\end{definition}

\begin{proposition}{Propriétés des probabilités}{}
    \begin{itemize}
        \item $P(\emptyset) = 0$
        \item Si $A, B \in \mathcal A$, $P(\bar A) = 1 - P(A)$
        \item Si $A \cap B = \emptyset$, $P(A \cup B) = P(A) + P(B)$
        \item Si $A \subset B$, $P(A) \leqslant P(B)$
        \item Pour tous $A, B \in \mathcal A$, $P(A \cup B) = P(A) + P(B) - P(A \cap B)$
    \end{itemize}
\end{proposition}

\begin{proof}
    Posons $P(\emptyset) = a \in [0, 1]$.

    $\forall n \in \N$, $A_n = \emptyset$ donc $\bigcup_{n \in \N} A_n = \emptyset$

    Donc par $\sigma$-additivité, $P(\emptyset) = P\left( \bigcup_{n \in \N} A_n \right) = \sum_{n \in \N} P(A_n) = \sum_{n \in \N} a$
    Donc $a = 0$.

    $A \in \mathcal A$, $A_0 = A$, $A_1 = \bar A$, $A_n = \emptyset$ pour $n \geqslant 2$.

    Par $\sigma$-additivité, $1 = P(\Omega) = P(A \cup \bar A) = P(A) + P(\bar A) $

    Si $A \cap B = \emptyset$, on pose $A_0 = A$, $A_1 = B$, $A_n = \emptyset$ pour $n \geqslant 2$.

    Donc $P(A \cup B) = P(A) + P(B)$ par $\sigma$-additivité.

    Si $A \subset B$, $B = A \cup (B \setminus A)$ et donc $P(B) = P(A) + P(B \setminus A) \geqslant P(A)$.

    $A \cup B = A \cup (B \setminus A\cap B)$ et $B = (A \cap B) \cup (B \setminus A \cap B)$

    Donc $P(A \cup B) = P(A) + P(B) + P(A \cap B)$ et $P(B) = P(A \cap B) + P(B \setminus A \cap B)$
\end{proof}


\subsubsection{Exemple fondamental : espaces probabilisés discrets }

\begin{definition}{Espace probabilisé discret}{}
    Soit $\Omega$ un ensemble non vide. On appelle distribution de probabilités 
    discrète sur $\Omega$ toute famille de $\R^+$ indexée par $\Omega$, sommable 
    et de somme égale à 1.     

    Une distribution de probabilités discrète sur $\Omega$ est donc une famille 
    $(p_\omega)_{\omega \in \Omega}$ telle que : 
    $$ \sum_{\omega \in \Omega } p_\omega = 1$$
\end{definition}

\begin{proposition}{Support d'une distribution de probas discrète}{}
    Une distribution de probas discrète $(p_\omega)_{\omega \in \Omega}$
    est de support au plus dénombrable, c'est-à-dire que 
    $\{\omega \in \Omega \mid p_\omega \neq 0 \}$ est au plus dénombrable.
\end{proposition}

\begin{proof}
    Le support d'une famille sommable est au plus dénombrable.
\end{proof}

\begin{theoreme}{Espace probabilisé discret}{}
    Soit $\Omega \neq \emptyset$ et $(p_\omega)_{\omega \in \Omega}$ une distribution de probabilités discrète
    sur $\Omega$. En posant $\mathcal A = \mathscr P(\Omega)$ et
    $\fonction{P}{\mathscr P(\Omega)}{[0, 1]}{A}{P(A) = \sum_{\omega \in A} p_\omega }$
    Alors $(\Omega, \mathcal A, P)$ est un espace probabilisé.

    $P$ est appelée probabilité associée à cette distribution. 
\end{theoreme}

\begin{remarque}
    On a en particulier, $\forall \omega \in \Omega, \, P(\{\omega\}) = p_\omega$
\end{remarque}

\begin{proof}
    $\mathcal A = \mathscr P(\Omega)$ est une tribu.

    $P$ est une proba : 

    Si $A \subset \Omega$, $P(A) = \sum_{\omega \in A} p_\omega \leqslant \sum_{\omega \in \Omega} p_\omega = 1$

    Donc $P(A) \in [0, 1]$.

    $\displaystyle P(\Omega) = \sum_{\omega \in \Omega} p_\omega = 1$

    Soit $(A_n)_{n \in \N}$ famille disjointe de parties de $\Omega$ et $A = \bigsqcup_{n \in \N} A_n$.

    Montrons la $\sigma$-additivité :
    $$ P(A) = \sum_{\omega \in A} p_\omega = \sum_{n \in \N} \sum_{\omega \in A_n} p_\omega = \sum_{n \in \N} P(A_n) $$
    par le théorème de sommation par paquets.

    Donc $P$ est une probabilité.
\end{proof}

\begin{theorement}{espace probabilisé discret}{}
    Soit $\Omega$ un ensemble non vide au plus dénombrable,
    $\mathcal A = \mathscr P(\Omega)$ et $P$ une probabilité sur $(\Omega, \mathcal A)$.

    Alors en posant $\forall \omega \in \Omega, \, p_\omega = P(\{\omega\})$,
    $(p_\omega)_{\omega \in \Omega}$ est une distribution de probabilités discrète 
    et $P$ est la probabilité associée à cette distribution.
\end{theorement}

\begin{proof}
    $\forall \omega, \, p_\omega \geqslant 0$ et $\displaystyle \Omega = \bigsqcup_{\omega \in \Omega} \{\omega\}$
    et on a $\displaystyle P(\Omega) = 1 = \sum_{\omega \in \Omega}  p_\omega$ donc 
    on a bien une distribution de probabilités discrète.

    De plus $\forall A \in \mathscr P(\Omega)$, 
    $\displaystyle A = \bigsqcup_{\omega \in A} \{\omega\}$ et donc $\displaystyle P(A) = \sum_{\omega \in A} P(\{\omega\}) = \sum_{\omega \in A} p_\omega$.

    Donc $P$ est bien la proba associée à la distribution. 
\end{proof}

\begin{remarque}
    Dans cette situation, pour $\omega \in \Omega, \, \{\omega\}$ est appelé un évènement élémentaire.
\end{remarque}

\begin{exemple}
    Soit $\Omega$ un ensemble fini non vide de cardinal $N$. Alors la famille 
    $(p_\omega)_{\omega \in \Omega}$ définie par
    $\forall \omega \in \Omega, \, p_\omega = \frac 1 N$
    est une distribution de probabilités discrète car 
    $0 \leqslant \frac 1 N $ et $\displaystyle \sum_{\omega \in \Omega} p_\omega = \sum_{\omega \in \Omega} \frac 1 N = N \cdot \frac 1 N = 1$.

    La proba $P$ associée est appelée probabilité uniforme sur $\Omega$ et vérifie 
    $$\forall A \in \mathscr P(\Omega), \, P(A) = \frac{\text{card}(A)}{N}$$
\end{exemple}

\begin{exemple}
    Jeu de pile ou face équilibré :
    $$ \Omega = \{0, 1\}$$
    $$ \mathcal A = \mathscr P(\Omega) = \{\emptyset, \{0\}, \{1\}, \{0, 1\}\}$$

    $P(\{0\}) = P(\{1\}) = \frac 1 2$

    Lancer de dé à 6 faces équilibré : 
    $ \Omega = \{1, 2, 3, 4, 5, 6\}$, $\mathcal A = \mathscr P(\Omega)$ et
    $\forall k \in \ninter 1 6, \, P(\{k\}) = \frac 1 6$. 
\end{exemple}

\subsection{Limites et continuités}

Soit $(\Omega, \mathcal A, P)$ un espace probabilisé.

\begin{theoreme}{Continuité croissante}{}
    Soit $(A_n)_{n \in \N}$ une suite croissante d'évènements, c'est à dire telle que $\forall n \in \N$, 
    $A_n \subset A_{n+1}$. Alors $$P(A_n) \xrightarrow[n \to + \infty]{} P\left( \bigcup_{n \in \N} A_n \right)$$
\end{theoreme}

\idee{
    Poser des évènements pour avoir une union disjointe. 
}

\begin{proof}
    Posons $A = \bigcup_{n \in \N} A_n$. $A$ est bien un évènement car $\mathcal A$ est une tribu.

    Posons $\forall k \in \N^*$, $B_k = A_k \setminus A_{k-1} = A_k \cap \bar A_{k-1} \in \mathcal A$
    on a $A_k = B_k \sqcup A_{k - 1}$. 

    Posons $B_0 = A_0$, donc $\forall n \in \N$, $\displaystyle A_n = \bigsqcup_{k = 0}^n B_k$ 
    et $A = \displaystyle \bigcup_{n \in \N} A_n = \bigsqcup_{k \in \N} B_k$.

    Or par $\sigma$-additivité, $\displaystyle P(A_n) = \sum_{k = 0}^n P(B_k)$ et 
    $\displaystyle P(A) = \sum_{k \in \N} P(B_k)$.

    Donc $P(A_n) \xrightarrow[n \to + \infty]{} P(A)$ par définition de la somme infinie.
\end{proof}

\begin{proposition}{}{}
    Soit $(c_n)_{n \in \N}$ une suite quelconque d'évènement. Alors 
    $$P\left( \bigcup_{k = 0}^n c_k \right) \xrightarrow[n \to + \infty]{} P\left( \bigcup_{k \in \N} c_k \right)$$
\end{proposition}

\idee{
    Utiliser la continuité croissante.
}

\begin{proof}
    On pose $\forall n \in \N, \, \displaystyle A_n = \bigcup_{k = 0}^n c_k$. 

    Alors $(A_n)_{n \in \N}$ est croissante, c'est à dire que $A_n \subset A_{n + 1}$. 

    Donc par continuité croissante, $P(A_n) \xrightarrow[n \to + \infty]{} P(A)$
    avec $\displaystyle A = \bigcup_{k \in \N} A_k = \bigcup_{k \in \N} c_k$.
\end{proof}

\begin{theoreme}{Continuité décroissante}{}
    Soit $(B_n)_{n \in \N}$ une suite décroissante d'évènements c'est à dire telle que $\forall n \in \N, \, B_{n + 1} \subset B_n$.
    Alors $$\displaystyle P(B_n) \xrightarrow[n \to + \infty]{} P\left( \bigcap_{n \in \N} B_n \right)$$
\end{theoreme}

\begin{proof}
    On pose $\forall n, \, A_n = \bar B_n = C_\omega B_n$ alors 
    $P(B_n) = 1 - P(A_n)$, or $A_n \subset A_{n + 1}$ donc par continuité croissante, 
    $\displaystyle P(A_n) \xrightarrow[n \to + \infty]{} P\left( \bigcup_{n \in \N} A_n \right) = P\left( \overline{ \bigcap_{n \in \N} B_n } \right)$.
\end{proof}

\begin{proposition}{}{}
    Soit $(C_n)_{n \in \N}$ une suite quelconque d'évènements. Alors
    $$P\left( \bigcap_{k = 0}^n C_k \right) \xrightarrow[n \to + \infty]{} P\left( \bigcap_{k \in \N} C_k \right)$$
\end{proposition}

\begin{proof}
    On applique la conséquence de la continuité croissante à $D_k = \bar C_k$.

    $\displaystyle P\left( \bigcup_{k = 0}^n D_k \right) \xrightarrow[n \to + \infty]{} P\left( \bigcup_{k \in \N} D_k \right)$

    Donc $\displaystyle 1 - P\left( \bigcap_{k = 0}^n C_k \right) \xrightarrow[n \to + \infty]{} 1 - P\left( \bigcap_{k \in \N} C_k \right)$
\end{proof}

\begin{theoreme}{Inégalité de Boole}{}
    Soit $(c_n)_{n \in \N}$ une suite quelconque  d'évènements. Alors 
    $$ P\left( \bigcup_{n \in \N} c_n \right) \leqslant \sum_{n \in \N} P(c_n) $$
    cette somme étant éventuellement égale à $+ \infty$.
\end{theoreme}

\begin{proof}
    Si les $c_k$ sont 2 à 2 disjoints, il y a égalité par $\sigma$-additivité.

    \uindent{Cas 1} 

    \svspace Si la série $(P(c_n))_{n \in \N}$ diverge, alors 
    $\displaystyle \sum_{k \in \N} P(c_k) = + \infty$ et l'inégalité est vérifiée.

    \uindent{Cas 2}
    
    \svspace Sinon, pour $n \in \N$, montrons par récurrence que 
    $$ P\left( \bigcup_{k = 0}^n c_k \right) \leqslant \sum_{k = 0}^n P(c_k) $$

    Hérédité : supposont que pour un certain $n \in \N$, $$P \underbrace{\left( \bigcup_{k = 0}^n c_k \right) }_{D_n} \leqslant \sum_{k = 0}^n P(c_k) $$

    Alors 
    \begin{align}
        P(D_{n + 1}) & = P(D_n \cup c_{n + 1}) \\
        & = P(D_n) + P(c_{n + 1}) - P(D_n \cap c_{n + 1}) \\
        & \leqslant P(D_n) + P(c_{n + 1}) \\
        & \leqslant \sum_{k = 0}^{n + 1} P(c_k)
    \end{align}

    En faisant tendre $n$ vers $+ \infty$ et en utilisant la continuité croissante, on obtient l'inégalité voulue.
\end{proof}

\subsection{Evènements négligeables et presque sûrs }

\begin{definition}{Évènements presque sûrs et négligeables}{}
    Soit $A \in \mathcal A$. On dit que $A$ est un évènement négligeable ssi $P(A) = 0$.

    On dit que $A$ est un évènement presque sûr ssi $P(A) = 1$.
\end{definition}

\begin{proposition}{}{stabilité des évènements négligeables et presque sûrs}{}
    Une réunion finie ou dénombrable d'évènements négligeables est négligeable. 

    Une intersection finie ou dénombrable d'évènements presque sûrs est presque sûre.
\end{proposition}

\idee{
    Utiliser l'inégalité de Boole.
}

\begin{proof}
    Soit $(A_n)_{n \in I}$ évènements négligeables avec $I$ fini ou dénombrable.

    Alors $\displaystyle P\left( \bigcup_{n \in I} A_n \right) \leqslant \sum_{n \in I} P(A_n) \leqslant 0$
    par inégalité de Boole.

    Donc $\displaystyle P\left( \bigcup_{n \in I} A_n \right) = 0$. 

    Soit $(B_n)_{n \in I}$ évènements presque sûrs. 
    \begin{align}
        P\left( \bigcap_{n \in I} B_n \right) & = 1 - P\left( \overline{ \bigcap_{n \in I} B_n } \right) \\
        & = 1 - P\left( \bigcup_{n \in I} \bar B_n \right) \\
        & \geqslant 1 - \sum_{n \in I} P(\bar B_n) \\
        & = 1 - 0 = 1
    \end{align}

\end{proof}

\begin{definition}{Système quasi-complet d'évènements}{}
    Soit $(A_n)_{n \in I}$ une famille d'évènements avec $I$ au plus dénombrable. On dit que cette famille
    finie est un système quasi-complet d'évènements ssi 
    \begin{itemize}
        \item Pour $i \neq j$, $A_i \cap A_j = \emptyset$
        \item $\displaystyle P\left( \bigcup_{n \in I} A_n \right) = 1$
    \end{itemize}

    c'est-à-dire que $\displaystyle \bigcup_{n \in I} A_n$ est un évènement presque sûr.
\end{definition}

\section{Probabilités conditionnelles et indépendance }

Soit $(\Omega, \mathcal A , P)$ un espace probabilisé.

\begin{theoreme}{Probabilité conditionnelle}{}
    Soit $B \in \mathcal A$ tel que $P(B) \neq 0$. Pour $A \in \mathcal A$, on pose 
    $$P_B(A) = P(A \mid B) = \frac{P(A \cap B)}{P(B)}$$

    alors $P_B$ est une probabilité sur $(\Omega, \mathcal A)$ appelée probabilité conditionnelle
    sachant $B$.
\end{theoreme}

\begin{proof}
    $P(B) \neq 0$ donc $P_B$ bien définie. 

    Soit $A \in \mathcal A$, $A \cap B \subset B$ donc $0 \leqslant P(A \cap B) \leqslant P(B)$

    Donc $0 \leqslant P_B(A) \leqslant 1$, donc $P_B$ est à valeurs dans $[0, 1]$.

    $P_B(\Omega) = \frac{P(\Omega \cap B)}{P(B)} = 1$ 

    $\sigma$-additivité : Soit $(A_n)_{n \in \N}$ suite d'évènements 2 à 2 disjoints.

    \begin{align}
        P_B\left( \bigsqcup_{n \in \N} A_n \right) & = \frac{P\left( \left( \bigsqcup_{n \in \N} A_n \right) \cap B \right)}{P(B)} \\
        & = \frac{P\left( \bigsqcup_{n \in \N} (A_n \cap B) \right)}{P(B)} \\
        & = \sum_{n = 0 }^{+ \infty } \frac{P(A_n \cap B)}{P(B)} \\
        & = \sum_{n = 0}^{+ \infty} P_B(A_n)
    \end{align}
\end{proof}

\subsection{Formule des probabilités somposées }

\begin{theoreme}{Formule des probabilités composées}{}
    Soient $A_1, \dots, A_n \in \mathcal A$ tel que $P(A_1 \cap \dots \cap A_{n - 1}) \neq 0$,
    alors $$P(A_1 \cap \dots A_n) = P(A_1) \cdot P(A_2 \mid A_1) \cdot P(A_3 \mid A_1 \cap A_2) \cdots P(A_n \mid A_1 \cap \dots \cap A_{n - 1})$$
\end{theoreme}

\begin{proof}
    $P(A_n \mid A_1 \cap \dots \cap A_{n - 1}) = \frac{P(A_1 \cap \dots \cap A_n)}{P(A_1 \cap \dots \cap A_{n - 1})} \times \cdots \times \frac{P(A_1 \cap A_2)}{P(A_1)} \times P(A_1) = P(A_1 \cap \dots \cap A_n)$
\end{proof}

\subsection{Formules des probabilités totales }

\begin{theoreme}{Formule des probabilités totales}{}
    Soit $I$ un ensemble au plus dénombrable et $(A_n)_{n \in I}$ un système quasi-complet d'évènements.
    Soit $B \in \mathcal A$, alors 
    $$P(B) = \sum_{n \in I} P(B \cap A_n)$$

    De plus, si $\forall n \in I, \, P(A_n) \neq 0$, alors
    $$P(B) = \sum_{n \in I} P(A_n) \cdot P(B \mid A_n)$$
\end{theoreme}

\begin{proof}
    Notons $\displaystyle C = \overline{ \bigsqcup_{n \in I} A_n }$

    A cause du système quasi-complet d'évènements, $P(C) = 0$ et 
    $\displaystyle \Omega = C \sqcup \left( \bigsqcup_{n \in I} A_n \right)$
    donc par $\sigma$-additivité,
    $$ P(B) = P(B \cap C) + \sum_{n \in I} P(B \cap A_n) $$
    or $B \cap C \subset C$ donc $P(B \cap C) = 0$, 
    donc $$ P(B) = \sum_{n \in I} P(B \cap A_n) $$

    Si de plus $P(A_n) \neq 0$, on a 
    $$ P(B \mid A_n) = \frac{P(B \cap A_n)}{P(A_n)} $$

    Donc $ P(B \cap A_n) = P(A_n) \cdot P(B \mid A_n) $
\end{proof}

\subsection{Formule de Bayes}

\begin{theoreme}{Formule de Bayes}{}
    Soient $A, B \in \mathcal A$ tels que $P(A) \neq 0$ et $P(B) \neq 0$.
    Alors 
    $$ P(A \mid B) = \frac{P(A) \cdot P(B \mid A)}{P(B)} $$
\end{theoreme}

\begin{proof}
    $P(A \cap B) = P(A \mid B) \cdot P(B) = P(B \mid A) \cdot P(A)$
\end{proof}

\begin{corollaire}{formule de Bayes généralisée}{}
    Soit $I$ au plus dénombrable, et soit $(A_i)_{i \in I}$ système quasi-complet d'évènements 
    de proba non nulles, et $B$ tel que $P(B) \neq 0$. Alors 
    $$ P(A_n \mid B) = \frac{P(A_n) \cdot P(B \mid A_n)}{\sum_{i \in I} P(A_i) \cdot P(B \mid A_i)} $$
\end{corollaire}

\subsection{Evènements indépendants}
\begin{definition}{Indépendance}{}
    Soient $A, B \in \mathcal A$. On dit que $A$ et $B$ sont indépendants ssi 
    $$ P(A \cap B) = P(A) \cdot P(B) $$
\end{definition}

\begin{proposition}{Caractérisation de l'indépendance}{}
    Soient $A, B \in \mathcal A$ tels que $P(B) \neq 0$. 

    Alors $A$ et $B$ sont indépendants ssi $P(A \mid B) = P(A)$.
\end{proposition}

\begin{proof}
    Si $A$ et $B$ sont indépendants, alors 
    $$ P(A \mid B) = \frac{P(A \cap B)}{P(B)} = \frac{P(A) \cdot P(B)}{P(B)} = P(A) $$

    Si $P(A \mid B) = P(A)$, alors
    $ \frac{P(A \cap B)  }{P(B) } = P(A) $ donc $P(A \cap B) = P(A) \cdot P(B)$
\end{proof}

\begin{proposition}{Indépendance avec le contraire}{}
    Si $A$ et $B$ sont indépendants, alors $A$ et $\overline B$ sont indépendants. 
\end{proposition}

\idee{
    Utiliser système complet d'évènements $\{B, \bar B\}$.
}

\begin{proof}
    Si $A$ et $B$ sont indépendants, alors 
    $P(A \cap \bar B) + P(A \cap B) = P(A)$
    donc $P(A \cap \bar B) = P(A) - P(A \cap B) = P(A) - P(A) \cdot P(B) = P(A) \cdot (1 - P(B)) = P(A) \cdot P(\bar B)$
\end{proof}

\begin{definition}{Indépendance mutuelle}{}
    Soit $I$ un ensemble non vide, et soit $(A_i)_{i \in I}$ une famille d'évènements.

    On dit que ces évènements sont mutuellement indépendants ssi pour tout
    ensemble fini $J \subset I$ non vide, on a 
    $$ P\left( \bigcap_{j \in J} A_j \right) = \prod_{j \in J} P(A_j) $$

    en particulier, si $j_1 \neq j_2 \in I$, alors 
    $P(A_{j_1} \cap A_{j_2}) = P(A_{j_1}) \cdot P(A_{j_2})$, donc 
    $A_{j_1}$ et $A_{j_2}$ sont indépendants. Les évènements sont donc deux à deux indépendants.
\end{definition}

\begin{proposition}{Lien indépendance mutuelle et deux à deux}{}
    L'indépendance d'évènements deux à deux n'implique pas l'indépendance mutuelle.
\end{proposition}

\begin{proof}
    On lance 2 dés équilibrés.

    A : le résultat du premier dé est pair.

    B : le résultat du second dé est pair.

    C : la somme des résultats est paire.

    $\Omega = \ninter 1 6 ^2$, $P(A) = P(B) = \frac 1 2$

    $P(A \cap B) = \frac 1 4 = P(A) \cdot P(B)$ car évènements indépendants.

    \begin{align}
        P(C) & = P(C \mid A) P(A) + P(C \mid \bar A) P(\bar A) \\
        & = P(B) \cdot P(A) + P(\bar B) \cdot P(\bar A) \\
        & = \frac 1 4 + \frac 1 4 = \frac 1 2 
    \end{align} 

    Donc $P(C \cap A) = P(B \cap A) = \frac 1 4 = P(C) \cdot P(A)$, donc 
    $A$ et $C$ sont indépendants.

    De même, $B$ et $C$ sont indépendants.

    $P(A \cap B \cap C) = P(A \cap B) = \frac 1 4 \neq P(A) \cdot P(B) \cdot P(C) = \frac 1 8$

    Donc $A, B, C$ sont deux à deux indépendants mais pas mutuellement indépendants.
\end{proof}

\section{Variables aléatoires }

Soit $(\Omega, \mathcal A, P)$ un espace probabilisé et $E \neq \emptyset$. 

\begin{definition}{Variable aléatoire discrète}{}
    On appelle variable aléatoire discrète sur $(\Omega, \mathcal A, P)$ à valeurs dans $E$
    toute application $\fonction{X}{\Omega}{E}{\omega}{X(\omega)}$ telle que 
    \begin{itemize}
        \item $X(\Omega)$ est au plus dénombrable
        \item $\forall x \in X(\Omega), \, X^{-1}(\{x\}) = \{\omega \in \Omega \mid X(\omega) = x\} \in \mathcal A$
    \end{itemize}

    On note alors l'évènement $\{X = x\} = X^{-1}(\{x\})$.
    
    Si $E = \R$, on dit que $X$ est une variable aléatoire réelle discrète.
\end{definition}

\begin{remarque}
    Si $\Omega$ est au plus dénombrable, et $\mathcal A = \mathscr P(\Omega)$, toute application 
    de $\Omega$ dans $E$ est une variable aléatoire discrète.
\end{remarque}

\subsection{Notations}

Ce paragraphe avant la def est potentiellement un peu bizare.

Soit $\fonction{X}{\Omega}{E}{\omega}{X(\omega)}$ une variable aléatoire discrète.

Soit $A \subset E$, alors $A \cap X(\Omega)$ est au plus dénombrable. 

Considérons $(x_i)_{i \in I}$ avec $I$ au plus dénombrable la famille des éléments 
de $A \cap X(\Omega)$, $A \cap X(\Omega) = \{x_i, \, i \in I\}$.

Ainsi $A \cap X(\Omega) = \bigsqcup_{i \in I} \{x_i\}$, or 
$X^{-1}({x_i}) = \{\omega \in \Omega, X(\omega) = x_i\} \in \mathcal A$, or $\mathcal A$
est une tribu, donc $\displaystyle \bigsqcup_{i \in I} X^{-1}(\{x_i\}) \in \mathcal A$.

Or $\bigsqcup_{i \in I} X^{-1}(\{x_i\}) = \{\omega \in \Omega, \exists i \in I, X(\omega = x_i)\} = \{\omega \in \Omega, X(\omega) \in A\}$


\begin{definition}{Évènements associés à une variable aléatoire}{}
    On notera alors $(X \in A) = X^{-1}(A) = \{\omega \in \Omega, X(\omega) \in A\}$ et donc 

    $$\displaystyle (X \in A) = \bigsqcup_{i \in I} (X = x_i)$$
    
\end{definition}

\begin{definition}{Notations pour les variables aléatoires réelles}{}
    Soit $\fonction{X}{\Omega}{\R}{\omega}{X(\omega)}$ une variable aléatoire réelle discrète.

    Alors pour tout $a \in \R$, on utilise les notations suivantes : 

    \begin{itemize}
        \item $(X \in [a, + \infty[) = (X \geqslant a)$
        \item $(X \in ]a, + \infty[) = (X > a)$
        \item $(X \in ]-\infty, a]) = (X \leqslant a)$
        \item $(X \in ]-\infty, a[) = (X < a)$
    \end{itemize}

    et par exemple, en notant $(x_i)_{i \in I}$ la famille des valeurs prises par $X$, 
    $\displaystyle (X \geqslant a) = \bigsqcup_{\substack{i \in I \\ x_i \geqslant a}} (X = x_i)$.
\end{definition}

\subsection{Loi d'une variable aléatoire discrète }

\begin{theoreme}{Loi d'une variable aléatoire discrète}{}
    Soit $X : \Omega \to E$ une variable aléatoire discrète. On considère l'espace probabilisable 
    $(X(\Omega), \mathscr P(X(\Omega)))$, et la fonction 
    $$\fonction{P_X}{\mathscr P(X(\Omega))}{[0, 1]}{A}{P(X^{-1}(A)) = P(X \in A)}$$

    Alors $P_X$ est une probabilité appelée loi de $X$. Cette loi est donc caractérisée par la donnée de : 
    \begin{itemize}
        \item L'ensemble $X(\Omega)$
        \item Pour tout $x \in X(\Omega)$, $P_X(\{x\}) = P(X = x)$
    \end{itemize}
    et on aura $\sum_{x \in X(\Omega)} P(X = x) = 1$

    De plus pour tout $A \subset E$, $\ds P(X \in A) = \sum_{x \in (X(\Omega) \cap A)} P(X = x)$
\end{theoreme}

\begin{proof}
    Si $A \in \mathscr P(X(\Omega))$, $X^{-1}(A) \in \mathcal A$ car $X$ est une variable aléatoire
    donc $P(X^{-1}(A))$ est bien définie et dans $[0, 1]$.

    \begin{align}
        P_X(X(\Omega)) & = P(X^{-1}(X(\Omega))) \\
        & = P(\Omega) = 1
    \end{align}

    \uindent{$\sigma$-additivité}

    \svspace Soit $(A_n)_{n \in \N}$ une suite de parties de $X(\Omega)$ 2 à 2 disjointes.

    \begin{align}
        P_X \left( \bigsqcup_{n \in \N} A_n \right) & = P\left( X^{-1}\left( \bigsqcup_{n \in \N} A_n \right) \right) \\
        & = P\left( \bigsqcup_{n \in \N} X^{-1}(A_n) \right) 
    \end{align}

    Soit $n_1 \neq n_2 \in \N$, montrons que $X^{-1}(A_{n_1}) \cap X^{-1}(A_{n_2}) = \emptyset$.

    Si $\omega \in X^{-1}(A_{n_1}) \cap X^{-1}(A_{n_2})$, alors $X(\omega) \in A_{n_1}$ et $X(\omega) \in A_{n_2}$,
    ce qui est impossible car $A_{n_1} \cap A_{n_2} = \emptyset$.

    Donc l'union est disjointe. 

    Donc $P_X(\bigsqcup_{n \in \N} A_n) = \sum_{n \in \N} P(X^{-1}(A_n)) = \sum_{n \in \N} P_X(A_n)$, 
    d'où la $\sigma$-additivité.

\end{proof}

\begin{definition}{Égalité de lois}{}
    Soit $X$ variable aléatoire discrète définie sur $(\Omega, \mathcal A, P)$, $Y$ variable aléatoire discrète définie 
    sur $(\Omega', \mathcal A', P')$ à valeurs dans $E$.

    On suppose que $X(\Omega) = Y(\Omega')$ et que $\forall x \in X(\Omega), \, P(X = x) = P'(Y = x)$.

    On dira alors que $X$ et $Y$ suivent la même loi, et on notera $X \sim Y$.
\end{definition}

\subsection{Composition}

\begin{theoreme}{Composition de variables aléatoires discrètes}{}
    Soit $F$ un ensemble non vide, $f : E \to F$ et $X : \Omega \to E$ une variable aléatoire discrète, 
    alors la composée $f \circ X : \Omega \to F$ est une variable aléatoire discrète notée $f(X)$.
\end{theoreme}

\begin{proof}
    $X(\Omega)$ est au plus dénombrable, donc $f(X(\Omega))$ est au plus dénombrable.

    Soit $y \in f(X(\Omega))$, 

    \begin{align}
        (f(X))^{-1}(\{y\}) & = \{\omega \in \Omega \mid f(X(\omega)) = y\} \\
        & = \{\omega \in \Omega \mid X(\omega) \in f^{-1}(\{y\})\} \\
        & \in \mathcal A
    \end{align}

    Donc $f(X)$ est une variable aléatoire discrète.
\end{proof}

\begin{theoreme}{}{}
    Soient $X, Y$ deux variables aléatoires discrètes telles que $X \sim Y$. 

    Soit $f : X(\Omega) \to F$, alors $f(X) \sim f(Y)$.
\end{theoreme}

\begin{proof}
    $X$ sur $(\Omega, \mathcal A, P)$, $Y$ sur $(\Omega', \mathcal A', P')$.

    On a $X \sim Y$, donc $X(\Omega) = Y(\Omega')$, donc $f(X(\Omega)) = f(Y(\Omega'))$.

    Soit $z \in f(X(\Omega))$, $P(f(X) = z) = P(X \in f^{-1}(\{z\})) = P'(Y \in f^{-1}(\{z\})) = P'(f(Y) = z)$.

    Donc $f(X) \sim f(Y)$.
\end{proof}

\subsection{Variables aléatoires indépendantes }

\begin{definition}{Indépendance de variables aléatoires discrètes}{}
    Soit $I \neq \emptyset$, et soit $(X_n)_{n \in I}$ des variables aléatoires discrètes. 

    On dit que ces variables aléatoires discrètes sont mutuellement indépendantes
    (ou indépendantes) ssi pour tout ensemble fini $J \subset I$, pour tous 
    $(x_j)_{j \in J}$ où $\forall n \in J, x_n \in X_n(\Omega)$, 
    $\displaystyle P\left( \bigcap_{j \in J} (X_j = x_j) \right) = \prod_{j \in J} P(X_j = x_j)$

    c'est à dire que les évènements $(X_j = x_j)$ sont mutuellement indépendants.
\end{definition}


\begin{proposition}{Indépendance des évènements associés}{}
    Soient $(X_n)_{n \in I}$ des variables aléatoires discrètes mutuellement indépendantes et 
    $(A_n)_{n \in I}$ une famille de parties des $X_n(\Omega)$. 

    Alors les évènements $(X_n \in A_n)_{n \in I}$ sont mutuellement indépendants, c'est à dire que pour 
    $J \subset I$ fini, $$ P\left( \bigcap_{j \in J} (X_j \in A_j) \right) = \prod_{j \in J} P(X_j \in A_j) $$
\end{proposition}

Preuve inutile 
 
\begin{proof}
    Pour $n \in J = \ninter 1 N$, notons 
    $(x_{i_n, n})_{i_n \in I_n}$ les éléments de $A_n$ famille au plus dénombrable, alors 
    $P(X_n \in A_n) = \sum_{i_n \in I_n} P(X_n = x_{i_n, n})$.

    Donc 
    \begin{align}
        \prod_{n \in J} P(X_n \in A_n) & = \sum_{i_1 \in I_1} \sum_{i_2 \in I_2} \cdots \sum_{i_N \in I_N} \prod_{n \in J} P(X_n = x_{i_n, n}) \\
        & = \sum \dots \sum P(\bigcap_{n \in J} (X_n = x_{i_n, n})) \\
        & = P(X_1 \in A_1 \cap \dots \cap X_n \in A_n)
    \end{align}
\end{proof}

\begin{theoreme}{Existence de suites de variables aléatoires indépendantes}{}
    Soit $(L_i)_{i \in \N}$ une suite de lois de variables aléatoires discrètes. Alors il existe 
    $(\Omega, \mathcal A, P)$ un espace probabilisé et $(X_i)_{i \in \N}$ une suite de variables aléatoires discrètes mutuellement indépendantes
    sur $\Omega$ telles que $\forall i \in \N, \, X_i \sim L_i$. 
\end{theoreme}

\begin{corollaire}{suite iid}{}
    Soit $L$ une loi de variables aléatoires discrètes. Il existe une suite 
    $(X_n)_{n \in \N}$ de variables aléatoires discrètes mutuellement indépendantes
    sur un espace probabilisé $(\Omega, \mathcal A, P)$ telles que $\forall n \in \N, \, X_n \sim L$.

    On dira alors que la suite de variables aléatoires discrètes $(X_n)_{n \in \N}$ 
    est iid, c'est à dire indépendante et identiquement distribuée.
\end{corollaire}

\subsection{Lois usuelles}

\begin{definition}{Loi uniforme}{}
    Soit $n \geqslant 1$ et $E$ un ensemble fini de cardinal $n$. Une variable aléatoire $X$ suit la 
    loi uniforme sur $E$ ssi 
    \begin{itemize}
        \item $X(\Omega) = E$
        \item $\forall x \in E, \, P(X = x) = \frac 1 n$
    \end{itemize}

    On notera alors $X \sim \mathcal U(E)$.
\end{definition}


\begin{definition}{Loi de Bernoulli}{}
    Soit $p \in ]0, 1[$. Une variable aléatoire discrète suit la loi de Bernoulli de paramètre 
    $p$ ssi 
    \begin{itemize}
        \item $X(\Omega) = \{0, 1\}$
        \item $P(X = 1) = p$ 
        \item $P(X = 0) = 1 - p$
    \end{itemize}
    On note alors $X \sim \mathcal B(p)$.
\end{definition}

\begin{definition}{Loi binomiale}{}
    Soit $p \in ]0, 1[$ et $n \geqslant 1$. Une variable aléatoire suit la loi binomiale de 
    paramètres $n$ et $p$ ssi 
    \begin{itemize}
        \item $X(\Omega) = \ninter 0 n$
        \item $\displaystyle \forall k \in \ninter 0 n, \, P(X = k) = \binom n k p^k (1 - p)^{n - k}$
    \end{itemize}
    On note alors $X \sim \mathcal B(n, p)$.
\end{definition}

\begin{remarque}
    Soient $X_1, \dots, X_n$ suite de variables aléatoires discrètes iid suivant $\mathcal B(p)$, 
    alors $\displaystyle X = \sum_{i = 1 }^{ n } X_i$ suit la loi $\mathcal B(n, p)$.
\end{remarque}

\begin{definition}{Loi géométrique}{}
    Soit $p \in ]0, 1[$. Une variable aléatoire $X$ suit la loi géométrique de paramètre $p$ ssi 
    \begin{itemize}
        \item $X(\Omega) = \N^*$
        \item $\forall k \in \N^*, \, P(X = k) = p(1 - p)^{k - 1} $
    \end{itemize}

    Cela définit bien une loi de variable aléatoire car $\displaystyle \sum_{k = 1}^{+ \infty} p(1 - p)^{k - 1} = 1$
    (formule somme géométrique).

    On notera alors $X \sim \mathcal G(p)$.
\end{definition}

\uindent{Interprétation }

\svspace Soit $(X_n)_{n \in \N^*}$ suite de variables aléatoires discrètes iid telles que $\forall n \in \N^*, \, X_n \sim \mathcal B(p)$.

Pour $\omega \in \Omega$, on note $X(\omega) = \begin{cases}
    + \infty & \text{ ssi } \forall n \in \N^*, X_n(\omega) = 0 \\
    k & \text{ ssi } X_k(\omega) = 1 \text{ et } \forall n < k, \, X_n(\omega) = 0
\end{cases}$

\begin{align}
    \displaystyle P(X = + \infty) & = P\left( \cap_{n \in \N} X_n = 0 \right) \\
    & = \lim_{N \to + \infty} P\left( \cap_{n = 1}^N X_n = 0 \right) \\
    & = \lim_{N \to + \infty} \prod_{n = 1 }^{N } P(X_n = 0) \\
    & = \lim_{N \to + \infty} (1 - p)^N = 0
\end{align}

$P(X = 1) = P(X_1 = 1) = p$

Si $k \geqslant 1$, 

\begin{align}
    P(X = k) & = P\left( X_1 = 0, \dots, X_{k - 1} = 0, X_k = 1 \right) \\
    & = p (1 - p)^{k - 1}
\end{align}

Donc $X \sim \mathcal G(p)$.

La loi géométrique est appelée loi du premier succès car elle modélise le nombre d'essais nécessaires pour obtenir le premier succès
dans une suite d'expériences de Bernoulli indépendantes.

\begin{definition}{Loi de poisson }{}
    Soit $\lambda > 0$. On dit qu'une variable aléatoire discrète $X$ suit la loi de 
    Poisson de paramètre $\lambda$ ssi 
    \begin{itemize}
        \item $X(\Omega) = \N$
        \item $\displaystyle \forall k \in \N, \, P(X = k) = \frac{e^{- \lambda} \lambda^k}{k!}$
    \end{itemize}

    On a bien $\displaystyle \sum_{k = 0}^{+ \infty} \frac{e^{- \lambda} \lambda^k}{k!} = e^{- \lambda} e^{\lambda} = 1$.

    On note alors $X \sim \mathcal P(\lambda)$.
\end{definition}

\uindent{Interprétation }

La loi de Poisson est la loi des évènements rares. Un évènement est dit rare lorsqu'il ne peut pas se produire 
plus d'une fois dans un laps de temps assez court. 

Soit $T > 0$ et $n \in \N^*$ assez grand pour que dans un intervalle de temps de durée 
$\frac T n$ l'évènement se produise au plus une fois. 

Donc dans chaque intervalle $\left[ \frac{kT} n, \frac{(k + 1)T}{n} \right[$ l'évènement se produit au plus
une fois, avec une proba $p = \alpha \frac T n$. 

Le nombre d'évènements dans l'intervalle $[0, T]$ suit donc une loi 
binomiale $\mathcal B\left( n, \alpha \frac T n \right)$.

Donc $\displaystyle \forall k \in \ninter 0 n$, $P(X = k) = \binom n k \left( \alpha \frac T n \right)^k \left( 1 - \alpha \frac T n \right)^{n - k}$.

Fixons $k$ et faisons tendre $n$ vers l'infini.

\begin{align}
    P(X = k ) & = \frac{n!}{k! (n - k)!} \frac{\alpha^k T^k}{n^k} e^{(n - k) \ln \left( 1 - \alpha \frac T n \right)} \\
    & \sim \frac{n^k}{k!} \frac{\alpha^k T^k}{n^k} e^{- \alpha T} \\
    & \xrightarrow[n \to + \infty]{} \frac{(\alpha T)^k}{k!} e^{- \alpha T}
\end{align}

Donc $X \sim \mathcal P(\alpha T)$. 

\section{Indépendance de variables aléatoires }

Soit $(\Omega, \mathcal A, P)$ un espace probabilisé.

\subsection{Loi conditionnelle}

\begin{definition}{Loi conditionnelle}{}
    Soit $X$ une variable aléatoire discrète et $A \in \mathcal A$ un évènement tel que $P(A) \neq 0$.

    On définit la loi conditionnelle de $X$ sachant $A$ par $\forall x \in X(\Omega)$,
    $$ P(X = x \mid A) = \frac{P(X = x \cap A)}{P(A)}$$
\end{definition}

\subsection{Couples de variables aléatoires discrètes }

\begin{theoreme}{Couple de variables aléatoires discrètes}{}
    Soient $X$ et $Y$ deux variables aléatoires discrètes sur $(\Omega, \mathcal A, P)$, on définit 
    leur couple $(X, Y)$ par l'application 
    $$\fonction{(X, Y)}{\Omega}{X(\Omega) \times Y(\Omega)}{\omega}{(X(\omega), Y(\omega))}$$
    alors c'est une variable aléatoire discrète.
\end{theoreme}

\begin{proof}
    $(X, Y)(\Omega) = \{(X(\omega), Y(\omega)), \omega \in \Omega \} \subset X(\Omega \times Y(\Omega))$ au plus dénombrable 
    donc $(X, Y)(\Omega)$ aussi. 

    Soit $(x, y) \in (X, Y)(\Omega)$. Alors 
    \begin{align}
        (X, Y)^{-1}(\{(x, y)\}) & = \{\omega \in \Omega, (X, Y)(\omega) = (x, y)\} \\
        & = \{\omega \in \Omega, X(\omega) = x \text{ et } Y(\omega) = y\} \\
        & = \{\omega \in \Omega, X(\omega) = x\} \cap \{\omega \in \Omega, Y(\omega) = y\} \\
        & = X^{-1}(\{x\}) \cap Y^{-1}(\{y\}) \in \mathcal A
    \end{align} 

    Donc $(X, Y)$ est une variable aléatoire discrète et $((X, y) = (x, y)) = (X = x) \cap (Y = y) = (X = x, Y = y)$.
\end{proof}

\begin{corollaire}{somme de variables aléatoires discrètes}{}
    Soient $X, Y$ deux variables aléatoires discrètes réelles sur $\Omega$, alors 
    $\fonction{X + Y}{\Omega}{\R}{\omega}{X(\omega) + Y(\omega)}$ est une variable aléatoire discrète 
    réelle. En effet, en posant $\fonction{f}{\R^2}{\R}{(x, y)}{x + y}$, on a $X + Y = f(X, Y)$ qui est 
    la composée de deux variables aléatoires discrètes, et donc en est une. 
\end{corollaire}

\begin{definition}{Lois conjointes et lois marginales}{}
    Soient $X, Y$ deux variables aléatoires discrètes sur $(\Omega, \mathcal A, P)$.
    La loi de la variable aléatoire discrète $(X, Y)$ est appelée loi conjointe. 

    Les lois de $X$ et $Y$ sont appelées lois marginales.
\end{definition}

\begin{proposition}{Lois marginales à partir de la loi conjointe}{}
    On peut calculer les lois marginales à partir de la loi conjointe : en effet, si 
    $x \in X(\Omega)$, 
    \begin{align}
        (X = x) & = \bigsqcup_{y \in Y(\Omega)} (X = x \cap Y = y) \\
        & = \bigsqcup_{y \in Y(\Omega)} ((X, Y) = (x, y))
    \end{align}

    Donc $\displaystyle P(X = x) = \sum_{y \in Y(\Omega)} P((X, Y) = (x, y))$.
    De même pour $Y$. 
\end{proposition}

\begin{exemple}
    Marche de l'ivrogne. 

    $X$ : déplacement horizontal.

    $Y$ : déplacement vertical.

    \uindent{Loi conjointe } 

    \svspace 
    
    \begin{itemize}
        \item $P((X, Y) = (1, 0)) 1 \frac 1 4$
        \item $P((X, Y) = (-1, 0)) 1 \frac 1 4$
        \item $P((X, Y) = (0, 1)) 1 \frac 1 4$
        \item $P((X, Y) = (0, -1)) 1 \frac 1 4$
    \end{itemize}

    Lois marginales : 
    \begin{itemize}
        \item $P(X = 1) = P((X, Y) = (1, 0)) = \frac 1 4$
        \item $P(X = -1) = P((X, Y) = (-1, 0)) = \frac 1 4$
        \item $P(X = 0) = P((X, Y) = (0, 1)) + P((X, Y) = (0, -1)) = \frac 1 2$
    \end{itemize}

    On peut représenter cela sous forme d'un tableau : 
    \begin{center}
        \begin{tabular}{c|c|c|c|c|c}
            $X \backslash Y$ & -1 & 0 & 1 & Total \\
            \hline
            -1 & 0 & $\frac 1 4$ & 0 & $\frac 1 4$ \\
            \hline
            0 & $\frac 1 4$ & 0 & $\frac 1 4$ & $\frac 1 2$ \\
            \hline
            1 & 0 & $\frac 1 4$ & 0 & $\frac 1 4$ \\
            \hline
            Total & $\frac 1 4$ & $\frac 1 2$ & $\frac 1 4$ & 1 
        \end{tabular}
    \end{center}


\end{exemple}

\begin{definition}{Indépendance de variables aléatoires discrètes}{}
    Deux variables aléatoires discrètes $(X, Y)$ sont indépendantes ssi 
    $\forall (x, y) \in X(\Omega) \times Y(\Omega), \, P((X, Y) = (x, y)) = P(X = x) \cdot P(Y = y)$.

    On notera alors $X \sqcup Y$.
\end{definition}

\begin{corollaire}{}{}
    En ce cas pour $A \subset X(\Omega)$, et $B \subset Y(\Omega)$,
    $P((X \in A) \cap (Y \in B)) = P(X \in A) \cdot P(Y \in B)$.
\end{corollaire}

\begin{proposition}{Indépendance par composition}{}
    Soit $X$ une variable aléatoire discrète à valeurs dans $E$, et 
    $Y$ une variable aléatoire discrète à valeurs dans $F$. 

    Soient $f : E \to E'$ et $g : F \to F'$. Alors si $X$ et $Y$ sont indépendantes,
    $f(X)$ et $g(Y)$ le sont également.
\end{proposition}

\begin{proof}
    Soient $a \in f(X(\Omega))$, et $b \in g(Y(\Omega))$. Donc $a \in E'$ et $b \in F'$.

    \begin{align}
        P(f(X) = a, g(Y) = b) & = P(X \in f^{-1}(\{a\}), Y \in g^{-1}(\{b\})) \\
        & = P(X \in f^{-1}(\{a\})) \cdot P(Y \in g^{-1}(\{b\})) \\
        & = P(f(X) = a) \cdot P(g(Y) = b)
    \end{align}

    Donc $f(X)$ et $g(Y)$ sont indépendantes.
\end{proof}

\subsection{Extension aux familles de n variables }

\begin{proposition}{N-uplet de variables aléatoires discrètes}{}
    Soient $(X_1, \dots, X_n)$ des variables aléatoires discrètes sur 
    $(\Omega, \mathcal A, P)$. On définit leur $n$-uplet 
    $\fonction{(X_1, \dots, X_n)}{\Omega}{X_1(\Omega) \times \cdots \times X_n(\Omega)}{\omega}{(X_1(\omega), \dots, X_n(\omega))}$.

    Il s'agit d'une variable aléatoire discrète. On dit que ces variables sont indépendantes 
    ssi $\forall k \leqslant n, \forall i_1, \dots, i_k \in \ninter 1 n, \forall (x_{i_1}, \dots, x_{i_k}) \in X_{i_1}(\Omega) \times \cdots \times X_{i_k}(\Omega)$,
    $$ P((X_{i_1}, \dots, X_{i_k}) = (x_{i_1}, \dots, x_{i_k})) = \prod_{j = 1}^k P(X_{i_j} = x_{i_j})$$
\end{proposition}

Les résultats du 2) se généralisent à ce cas. 

\subsection{Lemme des coalitions} 

\begin{theoreme}{Lemme des coalitions}{}
    Soient $X_1, \dots, X_n$ des variables aléatoires discrètes indépendantes. 
    Soit $m \in \ninter{1}{n - 1}$, et $f : X_1(\Omega) \times \cdots \times X_m(\Omega) \to E$,
    $g : X_{m + 1}(\Omega) \times \cdots \times X_n(\Omega) \to F$.

    Alors les variables aléatoires discrètes $f(X_1, \dots, X_m)$ et $g(X_{m + 1}, \dots, X_n)$ sont indépendantes.
\end{theoreme}

\begin{proof}
    Posons $Y = (X_1, \dots, X_m)$ et $Z = (X_{m + 1}, \dots, X_n)$.

    Les $(X_i)$ étant indépendants, $Y$ et $Z$ sont indépendants, donc par composition $f(Y)$ et $g(Z)$ sont indépendantes,
    c'est à dire $f(X_1, \dots, X_m)$ et $g(X_{m + 1}, \dots, X_n)$ sont indépendantes.
\end{proof}

\begin{proposition}{}{}
    Le lemme des coalitions reste vrai si $n$ considère plus de 2 coalitions. 
\end{proposition}

\section{Espérance}
Soit $(\Omega, \mathcal A, P)$ un espace probabilisé. 

\subsection{Variables aléatoires à valeurs positives}

\begin{definition}{Espérance d'une variable aléatoire positive}{}
    Soit $X$ une variable aléatoire discrète à valeurs dans $\R^+ \cup \{+\infty\}$.

    L'espérance de $X$ est définie comme étant la somme dans $[0, + \infty]$ 
    de la famille des $(x P(X = x) )_{x \in X(\Omega)}$. On la note $E(X)$.

    \uindent{Cas 1 } $P(X = + \infty) > 0$, alors $E(X) = + \infty$.

    \uindent{Cas 2 } $P(X = + \infty) = 0$, alors on a une famille de réels positifs. 
    Il y a alors deux cas : 
    \begin{itemize}
        \item Cette famille n'est pas sommable. Alors $E(X) = + \infty$.
        \item Cette famille est sommable. On dit que $X$ est d'espérance finie, et 
        $$ E(X) = \sum_{x \in X(\Omega)} x P(X = x) $$
    \end{itemize}
\end{definition}

\begin{remarque}
    Cette somme a bien un sens, car la famille est composée de réels positifs.
\end{remarque}

\begin{proposition}{Espérance vad sur N }{}
    Soit $X$ une variable aléatoire discrète à valeurs dans $\N \cup \{+\infty\}$. Alors
    $$ E(X) = \sum_{n = 1}^{+ \infty} P(X \geqslant n) $$
\end{proposition}

\begin{proof}
    Si $P(X = + \infty) >0$, $E(X) = + \infty$ et

    $\forall n \geqslant 1, \, P(X \geqslant n) \geqslant P(X = + \infty) > 0$, donc
    $\displaystyle \sum_{n = 1}^{+ \infty} P(X \geqslant n) \geqslant \sum_{n = 1}^{+ \infty} P(X = + \infty) = + \infty$.

    \avspace Sinon $\sum_{n \geqslant 1} P(X \geqslant n)$ est la somme d'une série à termes positifs
    convergente ou divergente.

    Donc 
    \begin{align}
        \sum_{n = 1}^{+ \infty} P(X \geqslant n) & = \sum_{n = 1}^{+ \infty} \sum_{k = n}^{+ \infty} P(X = k) \\
        & = \sum_{n = 1}^{+ \infty} \sum_{p = 0}^{+ \infty} \mathbb 1_{p \geqslant n} P(X = p) \\
        & = \sum_{p = 0}^{+ \infty} \sum_{n = 1}^{+ \infty} \mathbb 1_{p \geqslant n} P(X = p) \\
        & = \sum_{p = 0}^{+ \infty} p P(X = p) = E(X)
    \end{align}

    En utilisant le théorème de Fubini positif. 

\end{proof}

\begin{proposition}{Espérance loi uniforme}{}
    Soient $a, b \in \N$ et $X \sim \mathcal U(\ninter{a}{b})$. Alors
    $$ E(X) = \frac{a + b}{2} $$
\end{proposition}

\begin{proof}
    Variable aléatoire discrète à valeurs dans $\ninter{a}{b}$, donc l'espérance existe. 
    \begin{align}
        E(X) & = \sum_{n = a}^b n P(X = n) \\
        & = \sum_{n = a}^b \frac n {b - a + 1} \\
        & = \sum_{n = a }^{b } \frac 1 2 \frac{((n + 1)^2 - n^2 - 1)}{b - a + 1} \\
        & = \sum_{n = a }^{b } \left( \frac {(n + 1)^2 - n^2}{b - a + 1} \right) - \frac {b - a + 1} {2 (b - a + 1)} \\
        & = \frac{(b + 1)^2 - a^2}{2 (b - a + 1)} - \frac{b - a + 1}{2 (b - a + 1)} \\
        & = \frac{b^2 + 2b + 1 - a^2 - b + a - 1}{2 (b - a + 1)} \\
        & = \frac{b^2 - a^2 + b + a}{2 (b - a + 1)} \\
        & = \frac{(b - a)(b + a) + (b + a)}{2 ( b - a + 1)} \\
        & = \frac{(b - a + 1)(b + a)}{2 (b - a + 1)} \\
        & = \frac{a + b}{2}
    \end{align}

\end{proof}

\begin{proposition}{Espérance loi de Bernoulli}{}
    Soit $p \in ]0, 1[$ et $X \sim \mathcal B(p)$. Alors 
    $$ E(X) = p $$
\end{proposition}

\begin{proof}
    Si $X \sim \mathcal B(p)$, alors $X(\Omega) = \{0, 1\}$.
    Donc $E(X) = 0 \cdot P(X = 0) + 1 \cdot P(X = 1) = p$.
\end{proof}

\begin{proposition}{Espérance loi binomiale}{}
    Soit $n \geqslant 1$ et $p \in ]0, 1[$, et $X \sim \mathcal B(n, p)$. Alors 
    $$ E(X) = n p $$
\end{proposition}

\begin{proof}
    Si $X \sim \mathcal B(n, p)$, alors
       $$ E(X) = \sum_{k = 0}^n k \binom n k p^k (1 - p)^{n - k} $$
    
    or
    \begin{align}
        k \binom n k & = k \frac{n!}{k!(n - k)!} \\
        & = n \frac{n!}{(k - 1)! (n - 1 - (k - 1))!} \\
        & = n \frac{(n - 1)!}{(k - 1)!(n - 1 - (k - 1)!)} \\
        & = n \binom{n - 1}{k - 1}
    \end{align}

    Et donc 
    \begin{align}
        E(X) & = \sum_{k = 1}^n n \binom{n - 1}{k - 1} p^k (1 - p)^{n - k} \\
        & = n \sum_{k = 0 }^{n - 1} \binom{n - 1}{k} p^{k + 1} (1 - p)^{n - 1 - k} \\
        & = n p \sum_{k = 0}^{n - 1} \binom{n - 1}{k} p^k (1 - p)^{n - 1 - k} \\
        & = n p
    \end{align}
        
\end{proof}

\begin{proposition}{Espérance loi géométrique}{}
    Soit $p \in ]0, 1[$ et $X \sim \mathcal G(p)$. Alors 
    $$ E(X) = \frac 1 p $$
\end{proposition}

\begin{proof}
    Si $X \sim \mathcal G(p)$, alors 
    $$E(X) = \sum_{n = 1}^{+ \infty} n p (1 - p)^{n - 1}$$

    Pour $x \in ]-1, 1[$, posons $\displaystyle f(x) = \sum_{n = 0}^{+ \infty} x^n = \frac 1 {1 - x}$.

    Alors $f$ est somme de série entière de rayon de convergence $R = 1$, donc 
    $f$ est $\mathscr C^{1}$ sur $]-1, 1[$, et
    $\displaystyle f'(x) = \sum_{n = 1}^{+ \infty} n x^{n - 1} = \frac 1 {(1 - x)^2}$.

    $E(X) = p f'\left( 1 - p \right) = p \frac 1 {p^2} = \frac 1 p$.

\end{proof}

\begin{proposition}{Espérance loi de Poisson}{}
    Soit $\lambda > 0$ et $X \sim \mathcal P(\lambda)$. Alors 
    $$ E(X) = \lambda $$
\end{proposition}

\begin{proof}
    Si $X \sim \mathcal P(\lambda)$, alors 
    \begin{align}
        E(X) & = \sum_{k = 0}^{+ \infty} k \frac{e^{- \lambda} \lambda^k}{k!} \\
        & = e^{- \lambda} \sum_{k = 1}^{+ \infty} \frac{\lambda^k}{(k - 1)!} \\
        & = e^{- \lambda} \lambda \sum_{k = 0}^{+ \infty} \frac{\lambda^k}{k!} \\
        & = e^{- \lambda} \lambda e^{\lambda} = \lambda
    \end{align}
\end{proof}

\subsection{Variables aléatoires discrètes à valeurs complexes}

\begin{definition}{Espérance d'une variable aléatoire discrète complexe}{}
    Une variable aléatoire discrète dans $\C$ est dite d'espérance finie ssi 
    $(x(P(X = x)))_{x \in X(\Omega)}$ est sommable dans $\C$. On définit alors 
    son espérance par 
    $$ E(X) = \sum_{x \in X(\Omega)} x P(X = x) $$

    L'ensemble des variables aléatoires complexes d'espérance finie est notée $L^1$.
    Ainsi, $X \in L^1$ signifie que $E(X)$ existe dans $\C$, c'est à dire que 
    $(x P(X = x))_{x \in X(\Omega)}$ est sommable, c'est à dire que 
    $\displaystyle \sum_{x \in X(\Omega)} |x| P(X = x) < + \infty   $.
\end{definition}

\begin{definition}{Variables centrées}{}
    Soit $X \in L^1$. On dit que $X$ est centrée ssi 
    $$ E(X) = 0 $$    
\end{definition}

\subsection{Formule de transfert}

\begin{theoreme}{Formule de transfert}{}
    Soit $X$ une variable aléatoire discrète, et soit $f : X(\Omega) \to \C$.
    Alors $f \circ X$ que l'on note $f(X)$ est une variable aléatoire discrète complexe. 

    Alors $f(X)$ est d'espérance finie ssi la famille $(f(x) P(X = x))_{x \in X(\Omega)}$ 
    est d'espérance finie, et en ce cas
    $$ E(f(X)) = \sum_{x \in X(\Omega)} f(x) P(X = x) $$
\end{theoreme}

\begin{proof}
    Notons $X(\Omega) = \{x_i, i \in I\}$ ensemble au plus dénombrable. Alors 
    $f(X)(\Omega) = \{f(y_j), j \in J\}$ aussi.

    On a $X(\Omega) = \bigsqcup_{j \in J} \underbrace{f^{-1}(\{y_j\})}_{X_j}$.

    Considérons la famille $(f(x) P(X = x))_{x \in X(\Omega)}$ par le théorème de sommation 
    par paquets, elle est sommable ssi : 
    \begin{itemize}
        \item $\forall j \in J, \, (f(x) P(X = x))_{x \in X_j}$ est sommable
        \item La famille $\displaystyle \left(\sum_{x \in X_j} f(x) P(X = x)\right)_{j \in J}$ est sommable
    \end{itemize}

    Or si $x \in X_j$, $|f(x)| P(X = x) = |y_j| P(X = x)$ famille sommable 
    de somme $\displaystyle \sum_{x \in X_j} |f(x)| P(X = x) = |y_j| P(X \in X_j) = |y_j| P(f(X) = y_j)$.

    Donc la famille $(f(x) P(X = x))_{x \in X_j}$ est sommable ssi 

    $(|y| P(f(X) = y))_{y \in f(X(\Omega))}$ l'est 

    ssi $f(X)$ est d'espérance finie
    et en ce cas par théorème de sommation par paquets,
    \begin{align}
        \sum_{x \in X(\Omega)} f(x) P(X = x) & = \sum_{j \in J} y_j P(f(X) = y_j) \\
        & = \sum_{y \in f(X(\Omega))} y P(f(X) = y) \\
        & = E(f(X))
    \end{align}
\end{proof}

\subsection{Propriétés de l'espérance}

\begin{proposition}{Espace vectoriel et linéarité}{}
    $L^1$ est un $\C$-espace vectoriel, et l'espérance est linéaire, c'est à dire si 
    $X, Y \in L^1$, et $\alpha, \beta \in \C$, alors
    $$ E(\alpha X + \beta Y) = \alpha E(X) + \beta E(Y) $$
\end{proposition}

\idee{
    Commencer par majorer, puis calculer une fois l'existence établie.
}

\begin{proof}
    $L^1 \subset \mathscr F(\Omega, \C)$

    $o \in L^1$

    Soit $X \in L^1, \alpha \in \C^*$. 

    Considérons $\fonction{f}{X(\Omega)}{\C}{x}{\alpha x}$.
    Alors $f(X) \in L^1$ ssi $(\alpha x P(X = x))_{x \in X(\Omega)}$ est sommable,
    ssi $(x P(X = x))_{x \in X(\Omega)}$ est sommable, ssi $X \in L^1$.

    De plus par le théorème de transfert, 
    $E(f(X)) = E(\alpha X) = \sum_{x \in X(\Omega)} \alpha x P(X = x) = \alpha E(X)$.

    \lvspace Soient $X, Y \in L^1$.

    Considérons $Z = (X, Y)$ et $\fonction{f}{X(\Omega) \times Y(\Omega)}{\C}{(x, y)}{x + y}$.

    Alors par le théorème de transfert, $f(Z) = X + Y$ est d'espérance finie ssi

    $(f(x, y) p(Z = (x, y)))_{(x, y) \in X(\Omega) \times Y(\Omega)}$ est sommable.

    Et alors $\displaystyle E(X + Y) = \sum_{z \in f(X(\Omega) \times Y(\Omega))} f(z) P(Z = z)$.

    Considérons dans $\R \cup \{+ \infty \}$ 
    \begin{align}
        \sum_{(x, y) \in X(\Omega) \times Y(\Omega)} |f(x, y)| P(Z = (x, y)) & = \sum_{x \in X(\Omega)} \sum_{y \in Y(\Omega)} |x + y| P(X = x, Y = y) \\
        & \leqslant \sum_{x \in X(\Omega)} \sum_{y \in Y(\Omega)} (|x| + |y|) P(X = x, Y = y) \\
        & \leqslant \sum_{x \in X(\Omega)} |x| \sum_{y \in Y(\Omega)} P(X = x, Y = y) + \sum_{y \in Y(\Omega)} |y| \sum_{x \in X(\Omega)} P(X = x, Y = y) \\
        & \leqslant \sum_{x \in X(\Omega)} |x| P(X = x) + \sum_{y \in Y(\Omega)} |y| P(Y = y) \\
        & \leqslant E(|X|) + E(|Y|) < + \infty
    \end{align}

    car $X, Y \in L^1$.

    Et donc $X + Y$ est d'espérance finie, donc par le théorème de transfert, 
    \begin{align}
        E(X + Y) & = \sum_{(x, y) \in X(\Omega) \times Y(\Omega)} f(x, y) P(Z = (x, y)) \\
        & = E(X) + E(Y)
    \end{align}
    par le même calcul que précédemment.
\end{proof}



\begin{proposition}{Positivité de l'espérance}{}
    Soit $X \in L^1$ telle que $X \geqslant 0$, alors $E(X) \geqslant 0$.    
\end{proposition}

\begin{proof}
    Par définition, $E(X) = \sum_{x \in X(\Omega)} x P(X = x)$, 
    or $X \geqslant 0$ et donc $\forall x \in X(\Omega), \, x \geqslant 0$.
    Donc $E(X)$ est somme de réels positifs, donc $E(X) \geqslant 0$.
\end{proof}

\begin{proposition}{Croissance de l'espérance}{}
    Soient $X, Y \in L^1$ telles que $X \leqslant Y$, alors 
    $$ E(X) \leqslant E(Y) $$
    
\end{proposition}

\idee{
    Linéarité et positivité.
}

\begin{proof}
    $Y - X \in L^1$ par linéarité, et $Y - X \geqslant 0$.

    Donc par positivité de l'espérance, $E(Y - X) \geqslant 0$, donc 
    $E(Y) - E(X) \geqslant 0$, donc $E(X) \leqslant E(Y)$.
\end{proof}

\begin{proposition}{Inégalité triangulaire de l'espérance}{}
    Soit $X \in L^1$, alors $|X| \in L^1$ et
    $$ |E(X)| \leqslant E(|X|) $$
    
\end{proposition}

\idee{
    Formule de transfert et inégalité triangulaire classique.
}

\begin{proof}
    $X \in L^1$ ssi $(x P(X = x))_{x \in X(\Omega)}$ est sommable

    ssi $(|x| P(X = x))_{x \in X(\Omega)}$ est sommable

    ssi $|X| \in L^1$ par formule de transfert.

    Alors 
    \begin{align}
        |E(X)| & = \left| \sum_{x \in X(\Omega)} x P(X = x) \right| \\
        & \leqslant \sum_{x \in X(\Omega)} |x| P(X = x) \\
        & = E(|X|)
    \end{align} 
\end{proof}

\begin{proposition}{Critère de nullité espérance}{}
    Soit $X \in L^1$ telle que $X \geqslant 0$. Alors 
    $$ E(X) = 0 \iff P(X = 0) = 1 $$
    
\end{proposition}

\begin{proof}
    \begin{align}
        E(X) & = \sum_{x \in X(\Omega)} x P(X = x) \\
        & = 0
    \end{align}

    famille sommable de réels positifs, donc $\forall x \in X(\Omega), \, x P(X = x) = 0$, 

    donc $\forall x \in X(\Omega) \setminus \{0\}, \, P(X = x) = 0$.

    Or $\displaystyle 1 = \sum_{x \in X(\Omega)} P(X = x)$ donc $P(X = 0) = 1$.
\end{proof}

\begin{proposition}{Majoration et vad dans $L^1$}{}
    Soit $X, Y$ deux variables aléatoires discrètes complexes. Si 
    $|X| \leqslant Y$ et $Y \in L^1$, alors $X \in L^1$.
\end{proposition}


\begin{proof}
    $Y - |X| \geqslant 0$ variable aléatoire à valeurs dans $\R^+$ 
    
    donc $E(Y - |X|) \geqslant 0$ dans $\bar{\R}$, par linéarité $E(Y) - E(|X|) \geqslant 0$

    donc $E(|X|) \leqslant E(Y) < + \infty$.

    Donc $X \in L^1$.
\end{proof}

\subsection{Cas des variables aléatoires indépendantes }

\begin{theoreme}{Espérance du produit de variables indépendantes}{}
    Soient $X, Y \in L^1$ telles que $X \sqcup Y$. Alors 
    $XY \in L^1$ et $E(XY) = E(X) E(Y)$.
\end{theoreme}

\idee{
    Théorème de transfert et TGSCV.
}

\begin{proof}
    Posons $Z = (X, Y)$ et $\fonction{f}{X(\Omega) \times Y(\Omega)}{\C}{(x, y)}{xy}$.

    Donc $XY = f(Z)$.

    Par le théorème de transfert, $XY$ est d'espérance finie ssi $(f(x, y)P((X, Y) = (x, y)))_{(x, y) \in X(\Omega) \times Y(\Omega)}$ est sommable, 
    or pour $(x, y) \in X(\Omega) \times Y(\Omega)$, 
    \begin{align}
        f(x, y) P((X, Y) = (x, y)) & = x y P(X = x, Y = y) \\
        & = x P(X = x) \cdot y P(Y = y)
    \end{align}
    termes généraux de deux familles sommables, 
    
    donc $XY \in L^1$ et 
    $E(XY) = \displaystyle \sum_{x \in X(\Omega)} x P(X = x) \sum_{y \in Y(\Omega)} y P(Y = y) = E(X) E(Y)$.
\end{proof}

\begin{proposition}{}{}
    Soient $(X_1, \dots, X_p)$ dans $L^1$ mutuellement indépendantes, alors 
    $\displaystyle \prod_{j = 1}^p X_j \in L^1$ et 
    $$E\left( \prod_{j = 1}^{p } X_j \right) = \prod_{j = 1 }^{ p } E(X_j)$$
\end{proposition}

\idee{
    Récurrence et lemme des coalitions. 
}

\begin{proof}
    Par récurrence sur $p$. 

    Supposons le résultat vrai au rang $p$, et montrons le au rang $p + 1$.

    Soient $(X_1, \dots, X_{p + 1})$ dans $L^1$ mutuellement indépendantes.

    Posons $\displaystyle Z = \prod_{j = 1 }^{ p } X_j$, alors on a par HR 
    $Z \in L^1$ et $E(Z) = \prod_{j = 1 }^{ p } E(X_j)$.

    Par le lemme des coalitions, $Z$ et $X_{p + 1}$ sont indépendantes.

    Donc par le théorème précédent, 
    $\displaystyle \prod_{j = 1 }^{p + 1} X_j = Z X_{p + 1} \in L^1$ et
    $E(Z X_{p + 1}) = E(Z) E(X_{p + 1}) = \prod_{j = 1}^{p + 1} E(X_j)$.
\end{proof}

\section{Variance, covariance et écart-type }

Soit $(\Omega, \mathcal A, P)$ un espace probabilisé. Toutes les variables aléatoires 
considérées dans ce paragraphe sont réelles. 

\subsection{Variables aléatoires possédant un moment d'ordre 2 }

\begin{definition}{Moment d'ordre 2}{}
    Soit $X$ une variable aléatoire discrète réelle. 

    On dit que $X$ possède un moment d'ordre 2 ssi $X^2 \in L^1$, c'est à dire que
    $X^2$ est d'espérance finie.

    L'ensemble des variables aléatoires discrètes réelles possédant un moment d'ordre 2 
    est noté $L^2$.
\end{definition}



\begin{proposition}{Relation entre L2 et L1}{}
    Soit $X$ une variable aléatoire discrète réelle possédant un moment d'ordre $2$, 
    alors $X$ est d'espérance finie, c'est à dire $L^2 \subset L^1$.
\end{proposition}

\begin{proof}
    Soit $X \in L^2$. 

    On a $1 \in L^1$, or $(|X| -1 )^2 \geqslant 0$, donc 
    $X^2 - 2 |X| + 1 \leqslant 0$. 

    Donc $|X| \leqslant \frac 1 2[1 + X^2]$. 

    Donc $X \in L^1$. 
\end{proof}

\begin{proposition}{Nature de L2}{}
    $L^2$ est un $\R$-espace vectoriel. 
\end{proposition}

\uindent{Lemme } Si $X, Y \in L^2$, alors $XY \in L^1$.

\idee{
    Majorer $|XY|$ puis utiliser le fait que $L^2 \subset L^1$.
}

\begin{proof}
    $(|X| - |Y|)^2 \geqslant 0$, donc $|XY| \leqslant \frac 1 2 (X^2 + Y^2)$.

    Donc par majoration, $XY \in L^1$.

    $L^2 \subset L^1$ qui est un $\mathbb R$ espace vectoriel. 

    Soient $\alpha, \beta \in \R$ et $X, Y \in L^2$. Alors 

    \begin{align}
        (\alpha X + \beta Y)^2 & = \alpha^2 X^2 + \beta^2 Y^2 + 2 \alpha \beta XY \\
        & \leqslant \alpha^2 X^2 + \beta^2 Y^2 + 2 \alpha \beta XY
    \end{align}

    Donc par majoration, $(\alpha X + \beta Y)^2 \in L^1$, donc $\alpha X + \beta Y \in L^2$.
\end{proof}

\subsection{Inégalité de Cauchy-Schwarz }

\begin{theoreme}{Inégalité de Cauchy-Schwarz}{}
    Soient $X, Y \in L^2$, alors $XY \in L^1$ et
    $$ (E(XY))^2 \leqslant E(X^2) E(Y^2) $$
\end{theoreme}

\begin{proof}
    Pour tout réel $t \in \mathbb{R}$, considérons la variable aléatoire positive 
    
    $(tX + Y)^2$.
    Son espérance est donc positive ou nulle :
    $$ \forall t \in \mathbb{R}, \quad E\left( (tX + Y)^2 \right) \geqslant 0 $$
    
    En développant l'expression et en utilisant la linéarité de l'espérance, nous obtenons :
    $$ t^2 E(X^2) + 2t E(XY) + E(Y^2) \geqslant 0 $$
    
    Considérons cette expression comme un polynôme $P$ en la variable $t$, de la forme $P(t) = A t^2 + Bt + C$ avec :
    $$ A = E(X^2), \quad B = 2E(XY), \quad C = E(Y^2) $$
    
    Distinguons deux cas :
    
    \textbf{1. Cas principal :} Si $E(X^2) > 0$.
    $P(t)$ est un polynôme du second degré. Puisqu'il est positif ou nul pour tout $t$ (sa parabole est au-dessus de l'axe des abscisses), son discriminant $\Delta$ doit être négatif ou nul :
    $$ \Delta = B^2 - 4AC \leqslant 0 $$
    
    En remplaçant $A, B, C$ par leurs valeurs :
    $$ (2E(XY))^2 - 4 E(X^2) E(Y^2) \leqslant 0 $$
    $$ 4 (E(XY))^2 \leqslant 4 E(X^2) E(Y^2) $$
    
    En simplifiant par 4, nous obtenons l'inégalité de Cauchy-Schwarz :
    $$ (E(XY))^2 \leqslant E(X^2) E(Y^2) $$
    
    \textbf{2. Cas dégénéré :} Si $E(X^2) = 0$.
    Alors $X = 0$ presque sûrement, ce qui implique que le produit $XY = 0$ presque sûrement. 
    L'inégalité devient $0 \leqslant 0$, ce qui est toujours vérifié.
\end{proof}

\begin{proposition}{Cas d'égalité de Cauchy-Schwarz}{}
    Soient $X, Y \in L^2$, alors $E(XY)^2 = E(X^2) E(Y^2)$ ssi 
    $\begin{cases}
        P(X = 0) = 1 \\
        \text{ou} \\
        \exists \alpha \in \R, \, P(Y = \alpha X) = 1
    \end{cases}$
\end{proposition}

\begin{proof}
    Cas 1 : $E(X^2) = 0$ alors $E(XY) =0$ et égalité. En ce cas $P(X^2 = 0) = 1$ donc 
    $P(X = 0) = 1$.

    \uindent{Cas 2 } $E(X^2) \neq 0$ en ce cas égalité ssi $\Delta = 0$, c'est à dire
    
    ssi $\exists \alpha \in \R, \, E((\alpha X + Y)^2) = 0$

    ssi $\exists \alpha \in \R, \, P(\alpha X + Y = 0) = 1$
\end{proof}

\subsection{Variance et écart type }

\begin{definition}{Variance et écart type}{}
    Soit $X \in L^2$. On définit la variance de $X$ par : 
    $$V(X) = E\left( (X - E(X))^2 \right) $$

    Son écart type : 
    $$ \sigma(X) = \sqrt{V(X)} $$

    On dit que $X$ est réduite ssi $\sigma(X) = 1$. 
\end{definition}

\begin{proposition}{Critère de nullité de la variance}{}
    Soit $X \in L^2$, alors $V(X) = 0$ ssi 
    $\exists a \in \R$, $P(X = a) = 1$.
\end{proposition}

\begin{proof}
    Si $P(X = a) = 1$, alors $E(X) = a$.

    Donc $P((X - a)^2 = 0) = 1$, donc $E((X - a)^2) = 0$ donc $V(X) = 0$.

    \uindent{Réciproquement } Si $V(X) = 0$, posons $a = E(X)$.

    $E((X - a)^2) = 0$ or $(X - a)^2 \geqslant 0$, donc $P((X - a)^2 = 0) = 1$.
\end{proof}

\begin{proposition}{Affine de la variance}{}
    Soit $X \in L^2$, $a, b \in \R$. Alors $V(a X + b) = a^2 V(X)$ et 
    $\sigma(a X + b) = |a| \sigma(X)$.
\end{proposition}

\begin{proof}
    \begin{align}
        V(a X + b) & = E\left( (a X + b - E(a X + b))^2 \right) \\
        & = E\left( (a X - a E(X))^2 \right) \\
        & = a^2 E\left( (X - E(X))^2 \right) \\
        & = a^2 V(X)        
    \end{align}
\end{proof}

\begin{corollaire}{variable centrée réduite}{}
    Soit $X \in L^2$, tel que $\sigma \neq 0$. Alors 
    $Y = \frac{X - E(X)}{\sigma(X)}$ est une variable aléatoire centrée réduite.
\end{corollaire}

\begin{proof}
    $E(Y) = \frac{E(X) - E(X)}{\sigma(X)} = 0$ et 
    $\sigma(Y) = \frac{\sigma(X)}{\sigma(X)} = 1$.
\end{proof}

\begin{proposition}{Formule de König - Huygens}{}
    Soit $X \in L^2$, alors
    $$ V(X) = E(X^2) - (E(X))^2 $$
\end{proposition}

\begin{proof}
    \begin{align}
        V(X) & = E\left( (X - E(X))^2 \right) \\
        & = E(X^2 - 2 X E(X) + (E(X))^2) \\
        & = E(X^2) - 2 E(X) E(X) + (E(X))^2 \\
        & = E(X^2) - (E(X))^2
    \end{align}
\end{proof}

\subsubsection{Variance des lois classiques }

\begin{proposition}{Variance loi de Bernoulli}{}
    Soit $p \in ]0, 1[$ et $X \sim \mathcal B(p)$. Alors 
    $$ V(X) = p(1 - p) $$
\end{proposition}

\begin{proof}
    $X \sim \mathcal B(p)$ 

    $E(X^2) = 0^2 (1 - p) + 1^2 p = p$

    Donc $V(X) = E(X^2) - (E(X))^2 = p - p^2 = p(1 - p)$.
\end{proof}


\begin{proposition}{Variance loi binomiale}{}
    Soit $n \geqslant 1$ et $p \in ]0, 1[$, et $X \sim \mathcal B(n, p)$. Alors 
    $$ V(X) = n p (1 - p) $$
\end{proposition}

\idee{
    Utiliser la relation $E(X^2) = E(X(X - 1)) + E(X)$.
}

\begin{proof}
    $X \sim \mathcal B(n, p)$

    Pour $k \geqslant 1$, $\displaystyle k \binom n k = n \binom{n - 1}{k - 1}$

    $E(X^2) = E(X(X - 1) + X) = E(X(X - 1)) + E(X)$

    \begin{align}
        E(X(X - 1)) & = \sum_{k = 0}^n k (k - 1) \binom n k p^k (1 - p)^{n - k} \\
        & = \sum_{k = 2}^n n (k - 1) \binom {n - 1}{k - 1} p^k (1 - p)^{n - k} \\
        & = n \sum_{k = 2 }^{n } (n - 1) \binom{n - 2}{k - 2} p^k (1 - p)^{n - k} \\
        & = n (n - 1) p^2 \sum_{j = 0}^{n - 2} \binom{n - 2}{j} p^j (1 - p)^{n - 2 - j} \\
        & = n (n - 1) p^2
    \end{align}

    \begin{align}
        V(X) & = n(n - 1) p^2 + np - n^2 p^2 \\
        & = n p [(n - 1) p + 1 - n p] \\
        & = n p (1 - p)
    \end{align}
\end{proof}

\begin{proposition}{Variance loi géométrique}{}
    Soit $p \in ]0, 1[$ et $X \sim \mathcal G(p)$. Alors 
    $$ V(X) = \frac{1 - p}{p^2} $$
\end{proposition}

\begin{proof}
    Si $X \sim \mathcal G(p)$, alors
    $E(X(X - 1)) = \displaystyle \sum_{k = 2}^{+ \infty} k (k - 1) p (1 - p)^{k - 1}$ (transfert).

    Posons $\displaystyle f(x) = \sum_{k = 0}^{+ \infty} x^k = \frac 1 {1 - x}$.

    $x \in ]-1, 1[$, $f'(x) = \displaystyle \sum_{k = 1}^{+ \infty} k x^{k - 1} = \frac 1 {(1 - x)^2}$.

    $f''(x) = \displaystyle \sum_{k = 2}^{+ \infty} k (k - 1) x^{k - 2} = \frac 2 {(1 - x)^3}$.

    \begin{align}
        E(X(X - 1)) & = (1 - p) p f''(1 - p) \\
        & = (1 - p) p \frac 2 {p^3} \\
        & = \frac{2 (1 - p)}{p^2}
    \end{align}

    \begin{align}
        V(X) & = E(X(X - 1)) + E(X) - (E(X))^2 \\
        & = \frac{2 (1 - p)}{p^2} + \frac 1 p - \frac 1 {p^2} \\
        & = \frac{1 - p}{p^2}
    \end{align}
\end{proof}

\begin{proposition}{Variance et espérance loi de Poisson}{}
    Soit $\lambda > 0$ et $X \sim \mathcal P(\lambda)$. Alors 
    $$ V(X) = \lambda $$
\end{proposition}

\begin{proof}
    Si $X \sim \mathcal P(\lambda)$, alors 
    \begin{align}
        E(X(X - 1)) & = \sum_{k = 0}^{+ \infty} k (k - 1) \frac{e^{- \lambda} \lambda^k}{k!} \\
        & = e^{- \lambda} \sum_{k = 2}^{+ \infty} \frac{\lambda^k}{(k - 2)!} \\
        & = e^{- \lambda} \lambda^2 \sum_{k = 0}^{+ \infty} \frac{\lambda^k}{k!} \\
        & = e^{- \lambda} \lambda^2 e^{\lambda} = \lambda^2
    \end{align}

    \begin{align}
        V(X) & = E(X(X - 1)) + E(X) - (E(X))^2 \\
        & = \lambda^2 + \lambda - \lambda^2 = \lambda
    \end{align}
\end{proof}

\subsection{Covariance }

\begin{definition}{Covariance}{}
    Soient $X, Y \in L^2$. On définit leur covariance par : 
    $$ \cov(X, Y) = E\left( (X - E(X))(Y - E(Y)) \right) $$

    En particulier, $\cov(X, X) = V(X)$.
\end{definition}

\begin{proposition}{Formule de la covariance}{}
    Soient $X, Y \in L^2$. Alors 
    $$ \cov(X, Y) = E(XY) - E(X) E(Y) $$
\end{proposition}

\begin{proof}
    \begin{align}
        \cov(X, Y) & = E\left( (X - E(X))(Y - E(Y)) \right) \\
        & = E(XY - X E(Y) - Y E(X) + E(X) E(Y)) \\
        & = E(XY) - E(X) E(Y) - E(Y) E(X) + E(X) E(Y) \\
        & = E(XY) - E(X) E(Y)
    \end{align}
\end{proof}

\begin{definition}{Décorrélation}{}
    Soient $X, Y \in L^2$. On dit que $X$ et $Y$ sont \textbf{décorellées} ssi 
    $$ \cov(X, Y) = 0 $$
    c'est à dire ssi $E(XY) = E(X) E(Y)$.
\end{definition}

\begin{proposition}{Indépendance et décorellation}{}
    Si $X$ et $Y$ sont indépendantes, alors elles sont décorellées, donc $\cov(X, Y) = 0$.
\end{proposition}

\begin{remarque}
    La réciproque est fausse. 
\end{remarque}

\begin{theoreme}{Variance de la somme de variables}{}
    Soient $(X_i)_{1 \leqslant i \leqslant n}$ suite de $L^2$. Alors 
    $$ V\left( \sum_{i = 1}^n X_i \right) = \sum_{i = 1}^n V(X_i) + 2 \sum_{1 \leqslant i < j \leqslant n} \cov(X_i, X_j) $$
\end{theoreme}

\begin{corollaire}{variance de la somme de variables décorellées}{}
    Si les $X_i$ sont 2 à 2 décorellées, alors
    $$ V\left( \sum_{i = 1}^n X_i \right) = \sum_{i = 1}^n V(X_i) $$
\end{corollaire}

\begin{proof}
    \begin{align}
        V \left( \sum_{i = 1}^n X_i \right) & = E \left( \sum_{i = 1}^n X_i \right)^2 - \left( E \left( \sum_{i = 1}^n X_i \right) \right)^2 \\
        & = E \left( \sum_{i = 1}^n X_i^2 + 2 \sum_{1 \leqslant i < j \leqslant n} X_i X_j \right) - \left( \sum_{i = 1}^n E(X_i) \right)^2 \\
        & = \sum_{i = 1}^n E(X_i^2) + 2 \sum_{1 \leqslant i < j \leqslant n} E(X_i X_j) - \left( \sum_{i = 1}^n E(X_i)^2 + 2 \sum_{1 \leqslant i < j \leqslant n} E(X_i) E(X_j) \right) \\
    \end{align}
\end{proof}

\section{Inégalités probabilistes }
\subsection{Inégalité de Markov}

\begin{proposition}{Inégalité de Markov}{}
    Soit $X \in L^1$ et $a > 0$, alors 
    $$ P(|X| \geqslant a) \leqslant \frac{E(|X|)}{a} $$
\end{proposition}

\begin{proof}
    \begin{align}
        E(|X|) & = \sum_{x \in X(\Omega)} |x| P(X = x) \\
        & \geqslant \sum_{\substack{x \in X(\Omega) \\ |x| \geqslant a}} |x| P(X = x) \\
        & \geqslant \sum_{\substack{x \in X(\Omega) \\ |x| \geqslant a}} a P(X = x) \\
        & = a P(|X| \geqslant a)
    \end{align}
\end{proof}

\subsection{Inégalité de Bienaymé Tchebychev}

\begin{theoreme}{Inégalité de Bienaymé Tchebychev}{}
    Soit $X \in L^2$. Notons $m = E(X)$ et $\sigma = \sigma(X)$.
    Alors $$\forall \varepsilon > 0, P(|X - m| \geqslant \varepsilon) \leqslant \frac{\sigma^2}{\varepsilon^2}$$
\end{theoreme}

\begin{proof}
    Posons $Y = (X - m)^2$, $a = \varepsilon^2$.

    Par Markov, $P(|Y| \geqslant a) \leqslant \frac{E(|Y|)}{a}$.

    Donc $P(|X - m| \geqslant \varepsilon) \leqslant \frac{\sigma^2}{\varepsilon^2}$.
\end{proof}

\subsection{Loi faible des grands nombres }

\begin{proposition}{Une inégalité de concentration}{}
    Soit $(X_n)_{n \in \N}$ suite de variables aléatoires iid de $L^2$. 

    Notons $\displaystyle S_n = \sum_{k = 1 }^{n } X_k$ et $\forall n \in \N, \, m = E(X_n)$ et $\sigma = \sigma(X_n)$.
    Alors $$P\left( \left| \frac{S_n}{n} - m \right| \geqslant \varepsilon \right) \leqslant \frac{\sigma^2}{n \varepsilon^2}$$
\end{proposition}

\begin{proof}
    On pose $X = \frac{S_n}{n}$ et $\displaystyle E(X) = \frac{\sum_{i = 1 }^{n } E(X_i)} n$. 

    $\displaystyle V(X) = V\left( \frac 1 n \sum_{i = 1 }^{ n } X_i \right) = \frac 1 {n^2} \sum_{i = 1 }^{ n } V(X_i) = \frac{\sigma^2}{n}$.

    Par Bienaymé Tchebychev,
    $\displaystyle P\left( \left| \frac{S_n}{n} - m \right| \geqslant \varepsilon \right) = P(|X - E(X)| \geqslant \varepsilon) \leqslant \frac{\sigma^2}{n \varepsilon^2}$.
\end{proof}

\begin{proposition}{Loi faible des grands nombres}{}
    Soit $(X_n)_{n \in \N}$ suite iid de $L^2$. 
    $\displaystyle S_n = \sum_{k = 1 }^{ n } X_k$ et $m = E(X_n)$. Alors 
    $\forall \varepsilon > 0,$  $$P\left( \left| \frac{S_n}{n} - m \right| \geqslant \varepsilon \right) \xrightarrow[n \to + \infty]{} 0$$
\end{proposition}

\section{Fonctions génératrices }

Soit $(\Omega, \mathcal A, P)$ un espace probabilisé. 

\begin{definition}{Fonction génératrice}{}
    Soit $X$ une variable aléatoire à valeurs dans $\N$.

    On définit pour tout $t \in \R$, lorsque cette série converge, 
    $$ G_X(t) = \sum_{k = 0 }^{\infty } P(X = k) t^k $$
    appelée fonction génératrice de $X$.

    Lorsque $G_X(t)$ converge on a, par la formule de transfert,
    $$ G_X(t) = E(t^X) $$
\end{definition}

\begin{theoreme}{Rayon de convergence et continuité fonction génératrice}{}
    Soit $X$ une variable aléatoire à valeurs dans $\N$. La série entière 
    $G_X$ a un rayon de convergence supérieur ou égal à 1. 

    La convergence est normale sur $[-1, 1]$, donc $G_X$ est continue sur $[-1, 1]$.
\end{theoreme}

\begin{proof}
    En $t = 1$ la série converge donc le rayon de convergence $R \geqslant 1$.

    Soit $t $ tel que $|t| \leqslant 1$, alors 
    $|P(X = k) t^k| \leqslant P(X = k)$ terme général d'une série convergente d'où la 
    convergence normale, et la continuité sur $[-1, 1]$.
\end{proof}

\subsection{Cas des lois usuelles }

\begin{proposition}{Fonction génératrice loi de Bernoulli}{}
    Soit $p \in ]0, 1[$ et $X \sim \mathcal B(p)$. Alors 
    $$ G_X(t) = q + p t $$
    pour tout $t \in \R$, donc $R = + \infty$.
\end{proposition}

\begin{proof}
    Si $X \sim \mathcal B(p)$, $t \in [-1, 1]$, alors
    \begin{align}
        G_X(t) & = \sum_{k = 0}^1 P(X = k) t^k \\
        & = (1 - p) t^0 + p t^1 \\
        & = 1 - p + p t
    \end{align}
\end{proof}

\begin{proposition}{Fonction génératrice loi binomiale}{}
    Soit $n \geqslant 1$ et $p \in ]0, 1[$, et $X \sim \mathcal B(n, p)$. Alors 
    $$ G_X(t) = (1 - p + p t)^n $$
    pour tout $t \in \R$, donc $R = + \infty$.
\end{proposition}

\begin{proof}
    Si $X \sim \mathcal B(n, p)$, $t \in [-1, 1]$ et $n \in \N^*$, alors
    \begin{align}
        G_X(t) & = \sum_{k = 0}^n \binom n k p^k (1 - p)^{n - k} t^k \\
        & = \sum_{k = 0}^n \binom n k (p t)^k (1 - p)^{n - k} \\
        & = (1 - p + p t)^n
    \end{align}
\end{proof}

\begin{proposition}{Fonction génératrice loi géométrique}{}
    Soit $p \in ]0, 1[$ et $X \sim \mathcal G(p)$. Alors 
    $$ G_X(t) = \frac{pt}{1 - q t} $$
    pour tout $t \in \left]-\frac 1 q, \frac 1 q \right[$, donc 
    $R = \frac 1 q $.
\end{proposition}

\begin{proof}
    Si $X \sim \mathcal G(p)$, $t \in [-1, 1]$, alors
    \begin{align}
        G_X(t) & = \sum_{k = 1}^{\infty} p (1 - p)^{k - 1} t^k \\
        & = \frac{pt}{1 - (1 - p) t} \\
        & = \frac{pt}{1 - qt}
    \end{align}
\end{proof}

\begin{proposition}{Fonction génératrice loi de Poisson}{}
    Soit $\lambda > 0$ et $X \sim \mathcal P(\lambda)$. Alors 
    $$ G_X(t) = e^{\lambda (t - 1)} $$
    pour tout $t \in \R$, donc $R = + \infty$.
\end{proposition}

\begin{proof}
    Si $X \sim \mathcal P(\lambda)$, $t \in \R$, alors
    \begin{align}
        G_X(t) & = \sum_{k = 0}^{\infty} \frac{e^{- \lambda} \lambda^k}{k!} t^k \\
        & = e^{- \lambda} \sum_{k = 0}^{\infty} \frac{(\lambda t)^k}{k!} \\
        & = e^{- \lambda} e^{\lambda t} \\
        & = e^{\lambda (t - 1)}
    \end{align}
\end{proof}

\begin{theoreme}{Récupération de la loi à partir de la fonction génératrice}{}
    Soit $X$ une variable aléatoire à valeurs dans $\N$ de fonction génératrice $G_X$.
    Alors $$\forall k \in \N, \, P(X = k) = \frac{G_X^{(k)}(0)}{k!}$$
\end{theoreme}

\begin{corollaire}{égalité des lois à partir de la fonction génératrice}{}
    Soient $X, Y$ variables aléatoires discrètes à valeurs dans $\N$.
    Si $G_X = G_Y$, alors $X$ et $Y$ ont la même loi.
\end{corollaire}

\begin{proof}
    $G_X$ est un série entière de rayon de convergence non nul. 
\end{proof}

\subsection{Lien avec les moments}

\begin{theoreme}{}{}
    Soit $X$ une variable aléatoire à valeurs dans $\N$. Alors $X$ est 
    d'espérance finie ssi $G_X$ est dérivable en $1$ et dans ce cas
    $$ E(X) = G_X'(1) $$
\end{theoreme}

\begin{proof}
    $\Leftarrow$ : non exigible 

    $\Rightarrow$ : Supposons que $X$ est d'espérance finie.
    
    \uindent{Si $R = 1$ }

    Alors 
    $(k P(X = k))_{k \in \N}$ est sommable. 

    Or $\forall t \in [-1, 1]$, $G_X(t) = \sum_{k = 0}^{\infty} \underbrace{P(X = k) t^k}_{U_k(t)}$.

    $U_k$ est $\mathscr C^1$ et $u'_k(t) = k P(X = k) t^{k - 1}$.

    Donc $\forall t \in [-1, 1]$, $|u'_k(t)| \leqslant k P(X = k)$ terme général d'une série convergente 
    indépendante de $t$. 

    Donc $\sum u'_k$ converge normalement sur $[-1, 1]$.

    Donc $G_X$ dérivable en $1$ et $\displaystyle G'_X(1) = \sum_{k = 0}^{\infty} k P(X = k) 1^{k - 1} = E(X)$.

    \uindent{Si $R > 1$ }

    $G_X$ est dérivable en $1$ car $1$ est à l'intérieur de $]-R, R[$
    et $\forall t \in ]-R, R[$, $\displaystyle G'_X(t) = \sum_{k = 0}^{\infty} k P(X = k) t^{k - 1}$.
\end{proof}

\begin{corollaire}{variance avec fonction génératrice}{}
    Soit $X$ une variable aléatoire discrète à  valeurs dans $\N$ possédant un moment 
    d'ordre $2$, c'est à dire $X \in L^2$. Alors 
    $G_X$ est dérivable deux fois en $1$ et
    $$ G''_X(1) = E(X(X - 1)) $$

    Donc $$ V(X) = G''_X(1) + G'_X(1) - (G'_X(1))^2 $$
\end{corollaire}

\begin{proof}
    Si $X \in L^2$, si $R = 1$ alors 
    $\forall t \in [-1, 1]$ et 
    \begin{align}
        |u''_k(t)| & \leqslant k(k - 1) P(X = k) \\
        & \leqslant k^2 P(X = k)
    \end{align}
    terme général d'une série convergente indépendante de $t$.

    Donc $\sum_{k \geqslant 0} u''_k$ converge normalement sur $[-1, 1]$ et 
    $$G''_X(1) = \sum_{k = 0}^{\infty} k (k - 1) P(X = k) = E(X(X - 1))$$

    Si $R > 1$, $G_X$ est $\mathscr C^2$ sur $]-R, R[$ donc en dérivant 2 fois en 1 on obtient 
    la même chose.

    Donc \begin{align}
        V(X) & = E(X^2) - (E(X))^2 \\
        & = E(X(X - 1)) + E(X) - (E(X))^2 \\
        & = G''_X(1) + G'_X(1) - (G'_X(1))^2
    \end{align}
    car $X \in L^2$ donc $X \in L^1$.
\end{proof}

\subsection{Fonction génératrice d'une somme }

\begin{theoreme}{}{}
    Soient $X_1, \dots, X_n$ des variables aléatoires discrètes à valeurs dans $\N$ mutuellement 
    indépendantes. Alors $\forall t \in [-1, 1]$,
    $$ G_{X_1 + \cdots + X_n}(t) = \prod_{i = 1}^n G_{X_i}(t) $$
\end{theoreme}

\begin{proof}
    Soit $t \in [-1, 1]$.

    $G_{X_1 + \cdots + X_n}(t) = E\left( t^{X_1 + \cdots + X_n} \right)$

    Donc $G_{X_1 + \cdots + X_n}(t) = E\left( \prod_{i = 1 }^{n } t^{X_i} \right)$
    Or les $X_i$ sont mutuellement indépendantes, donc par le lemme des coalitions, 
    $(t^{X_i})_{1 \leqslant i \leqslant n}$ sont mutuellement indépendantes.

    Donc 
    \begin{align}
        G_{X_1 + \cdots + X_n}(t) & = \prod_{i = 1 }^{ n } E\left( t^{X_i} \right) \\
        & = \prod_{i = 1 }^{ n } G_{X_i}(t)
    \end{align}
\end{proof}
